% main.tex --- The Maya Book
%
% Build:   make            (in this directory; uses pdflatex twice)
% Output:  maya-book.pdf
%
% Chapter sources live under chapters/. To rearrange, edit \include lines.

\documentclass[11pt,openany,oneside]{book}

\input{preamble}

\date{}

\begin{document}

\frontmatter

% ----------------------------------------------------------------------
% Cover page --- one word, one line
% ----------------------------------------------------------------------
\begin{titlepage}
\thispagestyle{empty}

\vspace*{\fill}

\begin{center}
  {\fontsize{96}{100}\selectfont\sffamily\bfseries\color{mayablue} maya}
\end{center}

\vspace{0.6cm}

\begin{center}
  {\color{mayablue}\rule{4cm}{0.6pt}}
\end{center}

\vspace{0.6cm}

\begin{center}
  {\fontsize{12}{16}\selectfont\sffamily\color{mayagrey}
    a c\nolinebreak\hspace{-.05em}\raisebox{.2ex}{\small{++}}26 terminal ui library, from the byte up\par}
\end{center}

\vspace*{\fill}

\begin{center}
  {\fontsize{9}{12}\selectfont\sffamily\color{mayagrey}
    companion to the \textbf{agentty} sources\par}
\end{center}

\vspace*{1.5cm}
\end{titlepage}

% ----------------------------------------------------------------------
% Half-title / abstract
% ----------------------------------------------------------------------
\thispagestyle{empty}
\vspace*{\fill}
\begin{center}
\begin{minipage}{0.78\textwidth}
\small\itshape
This book is a deep guide to \maya{}, the header-mostly \cpp{}26
terminal-UI library that powers \agentty. We start with the terminal
as a machine, work up through the \cpp{} language features \maya{}
relies on, and then walk every layer of the library --- the DSL,
elements, the layout engine, the SIMD-accelerated renderer, the
signal graph, the Elm-style runtime, and the widget catalogue. Each
chapter ends with a working program.

\medskip

Everything in this book is grounded in the actual source under
\file{maya/}. Where the code disagrees with the prose, the code wins
--- report the mismatch as a bug.
\end{minipage}
\end{center}
\vspace*{\fill}
\clearpage

% ----------------------------------------------------------------------
% Preface
% ----------------------------------------------------------------------
\chapter*{Preface}
\addcontentsline{toc}{chapter}{Preface}

\section*{Who this book is for}

You have probably opened a terminal today. You have probably also
used a program that draws into it --- \code{htop}, \code{vim}, a git
client, a music player, a chat tool. This book is for the reader who
wants to write such programs themselves, in modern \cpp{}, and who
wants to understand every layer between the keystroke and the
flickering glyph.

It assumes nothing. You do not need to have written a TUI before. You
do not need to know \cpp{}26. You do not need to know what a pty is,
or what an ANSI escape sequence is, or why anyone would put a
compile-time string inside a type. Each of those is introduced from
zero, with the smallest example that makes the idea concrete, and
then connected to the place in \maya{}'s source where you will see
it again.

If you already know \cpp{}26 fluently and just want the \maya{}
tour, you can skim Chapter~\ref{ch:cpp}. Everyone else: read in
order. The book is a rope ladder; each rung sits on the one below
it.

\section*{What you will be able to do by the end}

Four concrete things, in this order:

\begin{enumerate}[leftmargin=*]
  \item Read the bytes flowing between a program and a terminal,
    and write a TUI ``Hello, red world'' in plain \cpp{} with no
    library at all (Chapter~\ref{ch:terminal}).
  \item Read every line of every \maya{} header --- templates,
    NTTPs, concepts, variants, \code{constexpr} chains --- and know
    what each construct is doing (Chapter~\ref{ch:cpp}).
  \item Build a non-trivial \maya{} program of your own, from a
    \code{Program} struct through a \code{view} function to a
    \code{Cmd<Msg>} effect, and trust that the compile-time checks
    catch the impossible states (Parts II and IV).
  \item Open the SIMD diff loop, the signal graph, the layout
    engine, or the inline-mode witness chain and follow what each
    one is doing on its own terms (Part III).
\end{enumerate}

That is the contract. If a chapter does not move you closer to one
of those four outcomes, file an issue.

\section*{How the book is shaped}

Four parts, mirroring four levels of zoom.

\paragraph{Part I --- Foundations.} The terminal as a state machine
speaking bytes; the \cpp{} language features \maya{} relies on. The
output of Part I is a reader who can run the examples and read the
headers.

\paragraph{Part II --- Architecture.} The Elm-style runtime
(\code{init}/\code{update}/\code{view}/\code{subscribe}), the
compile-time DSL that produces UI descriptions, the in-memory
\code{Element} tree, and the flexbox layout engine that places
them. The output of Part II is a reader who can build apps.

\paragraph{Part III --- Rendering.} Cells and canvases, SIMD-accelerated
frame diffing, the ANSI byte serialiser, and \maya{}'s unusual inline
mode that lives inside your scrollback instead of taking over the
terminal. The output of Part III is a reader who understands why
\maya{} feels fast.

\paragraph{Part IV --- Reactivity, runtime, and widgets.} The signal
graph that powers fine-grained reactivity, the input parser, the
full \code{Cmd}/\code{Sub} effect algebra, the catalogue of widgets,
and how to package and ship a binary built on \maya{}. The output
of Part IV is a reader who can ship.

\section*{How to read the boxes}

Five coloured boxes recur throughout. They are signposts; you can
read the book without them, but they tell you what the prose around
them is doing.

\begin{idea}
\textbf{Idea} (blue). One sentence summarising the key intuition of
the surrounding section. If you skim, read these.
\end{idea}

\begin{theory}
\textbf{Theory} (purple). The background concept from computer
science or programming-language theory behind what we just
explained. Skippable on first read; useful on second.
\end{theory}

\begin{warning}
\textbf{Pitfall} (amber). A mistake people new to the material
commonly make, and how to avoid it.
\end{warning}

\begin{recap}
\textbf{Recap} (green). A two-or-three-sentence summary of the
section just finished.
\end{recap}

\begin{terminalbox}
\textbf{Terminal} (black). Literal bytes or output, the way the
terminal sees or shows them.
\end{terminalbox}

Code in monospace with a coloured rule on the left is \cpp{} (or
occasionally shell). Lines are numbered. Long examples are split
into pieces and each piece is decoded after.

\section*{A note on accuracy}

Every code fragment in this book is either copied from \maya{}'s
headers or paraphrased from them. Where we simplify --- to keep an
example short, or to elide an implementation detail not yet
introduced --- we say so and name the file where the real version
lives. Read the book with the source open in another window;
that is the intended posture.

\bigskip

\noindent The chapters that follow assume nothing and explain
everything. Let us begin with the terminal.

\clearpage

\tableofcontents

\mainmatter

% ---- Part I: Foundations -----------------------------------------------
\part{Foundations}
\include{chapters/01-introduction}
\include{chapters/02-terminal}
\include{chapters/03-cpp}

% ---- Part II: Architecture ---------------------------------------------
\part{Architecture}
\include{chapters/04-elm}
\include{chapters/05-dsl}
\include{chapters/06-elements}
\include{chapters/07-layout}
\include{chapters/07a-concepts}
\include{chapters/07b-expected}
\include{chapters/07c-pipeline}
\include{chapters/07d-style}
\include{chapters/07e-builders}
\include{chapters/07f-text}
\include{chapters/07g-platform}
\include{chapters/07h-focus}
\include{chapters/07i-scroll}
\include{chapters/07j-errorboundary}
\include{chapters/07k-reactivity}
\include{chapters/07l-runtime}
\include{chapters/07m-testing}
\include{chapters/07n-widgetarch}
\include{chapters/07o-patterns}
\include{chapters/07p-performance}
\include{chapters/07q-guarantees}
\include{chapters/07r-cpp26}
\include{chapters/07s-comparison}
\include{chapters/07t-endtoend}
\include{chapters/07u-dispatch}
\include{chapters/07v-constexpr}
\include{chapters/07w-identity}
\include{chapters/07x-config}
\include{chapters/07y-observability}
\include{chapters/07z-i18n}
\include{chapters/07aa-lifecycle}
\include{chapters/07bb-a11y}
\include{chapters/07cc-security}
\include{chapters/07dd-cross}
\include{chapters/07ee-devx}
\include{chapters/07ff-summary}
\include{chapters/07gg-glossary}

% ---- Part III: Rendering -----------------------------------------------
\part{Rendering}
\include{chapters/08-canvas}
\include{chapters/09-simd}
\include{chapters/10-ansi}
\include{chapters/11-inline}
\include{chapters/11a-buffers}
\include{chapters/11b-diffloop}
\include{chapters/11c-colour}
\include{chapters/11d-width}
\include{chapters/11e-scheduler}
\include{chapters/11f-borders}
\include{chapters/11g-cursor}
\include{chapters/11h-mouse}
\include{chapters/11i-keyboard}
\include{chapters/11j-altscreen}
\include{chapters/11k-images}
\include{chapters/11l-hyperlinks}
\include{chapters/11m-titles}
\include{chapters/11n-clipping}
\include{chapters/11o-zorder}
\include{chapters/11p-pipeline2}
\include{chapters/11q-shaping}
\include{chapters/11r-allocations}
\include{chapters/11s-backend}
\include{chapters/11t-tracing}
\include{chapters/11u-snapshot}
\include{chapters/11v-bench}
\include{chapters/11w-slowlinks}
\include{chapters/11x-reconcile}
\include{chapters/11y-animations}
\include{chapters/11z-sprite}
\include{chapters/11ab-golden}
\include{chapters/11ac-reflow}

% ---- Part IV: Reactivity, runtime, widgets -----------------------------
\part{Reactivity, runtime, and widgets}
\include{chapters/12-signals}
\include{chapters/12a-signalgraph}
\include{chapters/12b-owners}
\include{chapters/12c-effects}
\include{chapters/12d-clocks}
\include{chapters/12e-async}
\include{chapters/12f-workers}
\include{chapters/12g-stores}
\include{chapters/12h-routing}
\include{chapters/12i-forms}
\include{chapters/12j-lists}
\include{chapters/12k-tables}
\include{chapters/12l-trees}
\include{chapters/12m-editors}
\include{chapters/12n-buttons}
\include{chapters/12o-progress}
\include{chapters/12p-menus}
\include{chapters/12q-dialogs}
\include{chapters/12r-toasts}
\include{chapters/12s-splits}
\include{chapters/12t-tabs}
\include{chapters/12u-statusbar}
\include{chapters/12v-pubsub}
\include{chapters/12w-fileio}
\include{chapters/12x-keymaps}
\include{chapters/12y-theming}
\include{chapters/12z-persistence}
\include{chapters/12aa-logging}
\include{chapters/12bb-plugins}
\include{chapters/12cc-pty}
\include{chapters/12dd-networking}
\include{chapters/12ee-a11y2}
\include{chapters/12ff-i18n2}
\include{chapters/12gg-security2}
\include{chapters/12hh-crossplat2}
\include{chapters/12ii-cmake}
\include{chapters/12jj-profile}
\include{chapters/12kk-memory}
\include{chapters/12ll-testing2}
\include{chapters/12mm-customwidget}
\include{chapters/12nn-curves}
\include{chapters/12oo-fsm}
\include{chapters/12pp-dnd}
\include{chapters/12qq-search}
\include{chapters/12rr-autocomplete}
\include{chapters/12ss-charts}
\include{chapters/12tt-pickers}
\include{chapters/12uu-diffview}
\include{chapters/12vv-markdown}
\include{chapters/12ww-syntax}
\include{chapters/12xx-undostate}
\include{chapters/12yy-hotreload}
\include{chapters/12zz-help}
\include{chapters/12ax-spinners}
\include{chapters/13a-offline}
\include{chapters/13b-selection}
\include{chapters/13c-signaldbg}
\include{chapters/13d-mvp}
\include{chapters/13e-lazy}
\include{chapters/13f-errorboundary2}
\include{chapters/13g-slots}
\include{chapters/13h-flexbox2}
\include{chapters/13i-grid}
\include{chapters/13j-optcases}
\include{chapters/13k-concurrency}
\include{chapters/13l-limits}
\include{chapters/13m-chatclient}
\include{chapters/13n-procmon}
\include{chapters/13o-filemgr}
\include{chapters/13p-musicplayer}
\include{chapters/13q-migrating}
\include{chapters/13r-distrib}
\include{chapters/13s-contributing}
\include{chapters/13t-catalogue1}
\include{chapters/13u-catalogue2}
\include{chapters/13v-catalogue3}
\include{chapters/13w-catalogue4}
\include{chapters/13x-catalogue5}
\include{chapters/13y-glossary2}
\include{chapters/13z-antipatterns}
\include{chapters/14a-futuregraphics}
\include{chapters/14b-futureinput}
\include{chapters/14c-futurelib}
\include{chapters/14d-web}
\include{chapters/14e-commonbugs}
\include{chapters/14f-recipes}
\include{chapters/14g-perfbudget}
\include{chapters/14h-loopannotated}
\include{chapters/14i-readingsource}
\include{chapters/14j-nextsteps}
\include{chapters/14k-colocation}
\include{chapters/14l-dataflow}
\include{chapters/14m-derived}
\include{chapters/14n-commands}
\include{chapters/14o-providers}
\include{chapters/14p-portals}
\include{chapters/14q-renderbudget}
\include{chapters/14r-apidesign}
\include{chapters/14s-docs}
\include{chapters/14t-v1}
\include{chapters/14u-conclusion}
\include{chapters/14v-uistate}
\include{chapters/14w-optimistic}
\include{chapters/14x-lazystate}
\include{chapters/14y-vscroll}
\include{chapters/14z-keydisc}
\include{chapters/15a-degradation}
\include{chapters/15b-progdisc}
\include{chapters/15c-empty}
\include{chapters/15d-formergo}
\include{chapters/15e-undobranches}
\include{chapters/15f-animchain}
\include{chapters/15g-timers}
\include{chapters/15h-flags}
\include{chapters/15i-exit}
\include{chapters/15j-afterword}
\include{chapters/15k-ansiref}
\include{chapters/15l-keyref}
\include{chapters/15m-signalref}
\include{chapters/15n-elementref}
\include{chapters/15o-troubleshoot}
\include{chapters/15p-further}
\include{chapters/15q-capref}
\include{chapters/15r-escglossary}
\include{chapters/15s-patternindex}
\include{chapters/15t-buildrecipes}
\include{chapters/15u-deployments}
\include{chapters/15v-partingchecklist}
\chapter{Reference manual: orientation}
\label{ch:refman}

\begin{aside}
The rest of this book taught you to think in \maya{}. This
appendix is the dictionary you reach for when you've already
internalised the grammar and just need to remember the exact
spelling of a verb.
\end{aside}

Every entry that follows obeys the same shape:

\begin{itemize}[leftmargin=*]
  \item A one-line gloss of what the thing is.
  \item The signature you'll type.
  \item A small but \emph{complete} program. Not a snippet —
    a file you can save, build, and run. Every example
    starts with the includes and ends with a \code{main()}
    or equivalent.
  \item A short \emph{notes} section for the rough edges.
\end{itemize}

\section{How to read these examples}

Every program in this manual compiles against the same
single header:

\begin{lstlisting}[language=C++]
#include <maya/maya.hpp>
using namespace maya;
using namespace maya::dsl;
\end{lstlisting}

Some widgets pull an extra header (\code{<maya/widget/...>}).
Those are listed inline where they appear.

The build line is always the same:

\begin{lstlisting}[language=bash]
g++ -std=c++23 example.cpp -lmaya -o example
\end{lstlisting}

\noindent or, with CMake:

\begin{lstlisting}[language=cmake]
find_package(maya REQUIRED)
add_executable(example example.cpp)
target_link_libraries(example PRIVATE maya::maya)
\end{lstlisting}

\section{Layout of this appendix}

Twenty-eight chapters. Each is a slice of the surface area:

\begin{itemize}[leftmargin=*]
  \item Programs and the runtime loop.
  \item Elements: text, layout, decoration.
  \item Style: colour, weight, modifiers.
  \item Subscriptions: keys, timers, signals.
  \item Commands: side effects as data.
  \item Widgets: every shipped widget, with a working demo.
  \item Inline rendering: \code{print} and \code{live}.
  \item Edge cases: resize, exit, errors, focus.
\end{itemize}

\section{The example template}

Use this as a starting point when you copy from this manual:

\begin{lstlisting}[language=C++]
#include <maya/maya.hpp>

using namespace maya;
using namespace maya::dsl;

struct App {
    struct Model {};
    struct Quit {};
    using Msg = std::variant<Quit>;

    static Model init() { return {}; }

    static auto update(Model m, Msg)
        -> std::pair<Model, Cmd<Msg>>
    {
        return {m, Cmd<Msg>::quit()};
    }

    static Element view(const Model&) {
        return t<"hello"> | Bold;
    }

    static auto subscribe(const Model&) -> Sub<Msg> {
        return key_map<Msg>({ {'q', Quit{}} });
    }
};

int main() { run<App>({.title = "demo"}); }
\end{lstlisting}

\noindent Every working example in this appendix is a
variation on that skeleton. If you forget where to start,
start here.

\section{A note on the prose}

These chapters are deliberately terser than the rest of the
book. You're here for the API, not the philosophy. Bullets
where bullets work; code where code works; one paragraph
of context per chapter, then on to the next entry.

\section{Recap}

\begin{recap}
\begin{itemize}[leftmargin=*]
  \item Every example is a complete program.
  \item Every program follows the same skeleton.
  \item Twenty-eight chapters cover every public surface.
  \item Read on demand, not in order.
\end{itemize}
\end{recap}

\section*{Workshop}

\begin{enumerate}[leftmargin=*]
  \item Type the template above into a file and build it.
  \item Run it; press \code{q}.
  \item You now have a working \maya{} project. The rest of
    this appendix is just things you can substitute in.
\end{enumerate}

\chapter{Reference: the Program concept}
\label{ch:refman-program}

\begin{aside}
The smallest contract a \maya{} app must satisfy. If your
struct passes \code{static\_assert(Program<App>)}, it will
run.
\end{aside}

\section{Required members}

A type \code{P} satisfies \code{Program} when it provides:

\begin{itemize}[leftmargin=*]
  \item a nested \code{Model} type;
  \item a nested \code{Msg} type (usually a
    \code{std::variant});
  \item \code{static Model init()};
  \item \code{static std::pair<Model, Cmd<Msg>>
    update(Model, Msg)};
  \item \code{static Element view(const Model\&)};
  \item \code{static Sub<Msg> subscribe(const Model\&)}.
\end{itemize}

\section{Minimal program}

\begin{lstlisting}[language=C++]
#include <maya/maya.hpp>

using namespace maya;
using namespace maya::dsl;

struct Hello {
    struct Model {};
    struct Quit {};
    using Msg = std::variant<Quit>;

    static Model init() { return {}; }

    static auto update(Model m, Msg)
        -> std::pair<Model, Cmd<Msg>>
    {
        return {m, Cmd<Msg>::quit()};
    }

    static Element view(const Model&) {
        return t<"press q to quit"> | Bold;
    }

    static auto subscribe(const Model&) -> Sub<Msg> {
        return key_map<Msg>({ {'q', Quit{}} });
    }
};

static_assert(Program<Hello>);

int main() {
    run<Hello>({.title = "hello"});
}
\end{lstlisting}

\section{Notes}

\begin{itemize}[leftmargin=*]
  \item \code{static\_assert(Program<Hello>)} is optional but
    cheap. It surfaces missing members at compile time
    rather than as deeply nested template errors.
  \item \code{init()} can take config args if you wrap
    \code{run<>} accordingly; the bare form takes none.
  \item \code{update} returns a \emph{pair}, not just a
    Model. The second slot is the command list.
  \item Use \code{Cmd<Msg>\{\}} for "no side effect."
\end{itemize}

\section{Common mistakes}

\begin{itemize}[leftmargin=*]
  \item Returning \code{Model} from \code{update} instead of
    a pair. The concept fails; the error is loud.
  \item Forgetting \code{Quit} in the variant. Your program
    runs but cannot exit cleanly.
  \item Putting state in static variables instead of the
    Model. Works once; breaks the moment you want to test.
\end{itemize}

\section{Recap}

\begin{recap}
\begin{itemize}[leftmargin=*]
  \item Six members; one concept.
  \item Pure functions; data effects.
  \item Verify with a static assert.
\end{itemize}
\end{recap}

\section*{Workshop}

\begin{enumerate}[leftmargin=*]
  \item Delete \code{subscribe} from the example. What does
    the compiler say?
  \item Restore it; delete \code{init}. Same question.
  \item Add a second variant case and a key for it.
\end{enumerate}

\chapter{Reference: \code{run} and \code{RunConfig}}
\label{ch:refman-run}

\begin{aside}
\code{run<P>} is the front door. It hands control of stdin,
stdout, and the terminal modes to \maya{}, executes your
program, then politely puts everything back.
\end{aside}

\section{Signature}

\begin{lstlisting}[language=C++]
template <Program P>
int run(RunConfig cfg = {});
\end{lstlisting}

\noindent Returns the process exit code (\code{0} on
graceful quit).

\section{RunConfig fields}

\begin{itemize}[leftmargin=*]
  \item \code{title}: window title, where supported.
  \item \code{alt\_screen}: switch to the alternate screen
    on start, restore on exit. Default \code{true}.
  \item \code{mouse}: enable mouse reporting. Default
    \code{false}.
  \item \code{fps}: target frame rate. Default \code{60}.
  \item \code{cursor}: \code{Hidden}, \code{Block}, or
    \code{Bar}. Default \code{Hidden}.
\end{itemize}

\section{Working example}

\begin{lstlisting}[language=C++]
#include <maya/maya.hpp>

using namespace maya;
using namespace maya::dsl;

struct Demo {
    struct Model { int frames = 0; };
    struct Tick {};
    struct Quit {};
    using Msg = std::variant<Tick, Quit>;

    static Model init() { return {}; }

    static auto update(Model m, Msg msg)
        -> std::pair<Model, Cmd<Msg>>
    {
        return std::visit(overload{
            [&](Tick) { m.frames++; return std::pair{m, Cmd<Msg>{}}; },
            [](Quit)  { return std::pair{Model{}, Cmd<Msg>::quit()}; },
        }, msg);
    }

    static Element view(const Model& m) {
        return v(
            t<"run() demo"> | Bold,
            blank_,
            text(std::format("frames: {}", m.frames)),
            blank_,
            t<"q to quit"> | Dim
        ) | pad<1> | border_<Round>;
    }

    static auto subscribe(const Model&) -> Sub<Msg> {
        return Sub<Msg>::batch(
            key_map<Msg>({ {'q', Quit{}} }),
            Sub<Msg>::every(std::chrono::milliseconds(16), Tick{})
        );
    }
};

int main() {
    return run<Demo>({
        .title      = "demo",
        .alt_screen = true,
        .mouse      = false,
        .fps        = 60,
    });
}
\end{lstlisting}

\section{Notes}

\begin{itemize}[leftmargin=*]
  \item Returning from \code{main} after \code{run} is fine;
    the runtime restores the terminal on its own.
  \item If \code{alt\_screen} is \code{false}, output stays
    in scrollback. Good for tools whose output you want to
    \code{grep} later.
  \item \code{fps} caps the loop, it does not force frames.
    Idle frames are skipped.
  \item Setting \code{mouse = true} is necessary before any
    mouse event will reach your \code{subscribe}.
\end{itemize}

\section{Recap}

\begin{recap}
\begin{itemize}[leftmargin=*]
  \item One function, one config struct.
  \item Title, alt-screen, mouse, fps, cursor.
  \item Returns an exit code.
\end{itemize}
\end{recap}

\section*{Workshop}

\begin{enumerate}[leftmargin=*]
  \item Flip \code{alt\_screen} to \code{false}. Watch the
    UI render in scrollback.
  \item Lower \code{fps} to \code{5}. Notice the change.
  \item Enable \code{mouse} and add a mouse handler in
    \code{subscribe}.
\end{enumerate}

\chapter{Reference: text elements}
\label{ch:refman-text}

\begin{aside}
\code{text} for runtime strings; \code{t<"...">} for literals
known at compile time. Both produce \code{Element}.
\end{aside}

\section{Constructors}

\begin{itemize}[leftmargin=*]
  \item \code{text(std::string)} — owns the bytes.
  \item \code{text(std::string\_view)} — borrows; caller
    must outlive the frame.
  \item \code{text(int)}, \code{text(double)} — formatted as
    decimal.
  \item \code{t<"literal">} — zero-allocation form for
    compile-time constants.
\end{itemize}

\section{Working example}

\begin{lstlisting}[language=C++]
#include <maya/maya.hpp>

using namespace maya;
using namespace maya::dsl;

struct TextDemo {
    struct Model { int n = 0; std::string label = "live"; };
    struct Bump {}; struct Quit {};
    using Msg = std::variant<Bump, Quit>;

    static Model init() { return {}; }

    static auto update(Model m, Msg msg)
        -> std::pair<Model, Cmd<Msg>>
    {
        return std::visit(overload{
            [&](Bump) { m.n++; return std::pair{m, Cmd<Msg>{}}; },
            [](Quit)  { return std::pair{Model{}, Cmd<Msg>::quit()}; },
        }, msg);
    }

    static Element view(const Model& m) {
        return v(
            t<"compile-time literal">,
            text("runtime owned string"),
            text(m.label),                    // dynamic
            text(m.n),                        // formatted int
            text(std::format("{:.2f}", 3.14)) // your own format
        ) | pad<1> | border_<Round>;
    }

    static auto subscribe(const Model&) -> Sub<Msg> {
        return key_map<Msg>({
            {'q', Quit{}},
            {' ', Bump{}},
        });
    }
};

int main() { run<TextDemo>({.title = "text"}); }
\end{lstlisting}

\section{Notes}

\begin{itemize}[leftmargin=*]
  \item \code{t<"...">} is preferred whenever the string is
    a literal: it skips the allocation and the copy.
  \item For format strings, build the \code{std::string}
    first; pass the result to \code{text}.
  \item Unicode is supported; the layout engine treats wide
    glyphs as two cells.
\end{itemize}

\section{Recap}

\begin{recap}
\begin{itemize}[leftmargin=*]
  \item \code{t<>} for literals, \code{text} for runtime.
  \item Numeric overloads exist for convenience.
  \item Use \code{std::format} for anything fancy.
\end{itemize}
\end{recap}

\section*{Workshop}

\begin{enumerate}[leftmargin=*]
  \item Replace one \code{text(...)} with \code{t<"...">}
    where the value is a literal.
  \item Show the current time using
    \code{std::chrono::system\_clock}.
  \item Render the value of a \code{std::vector<int>}'s size
    via \code{text(int)}.
\end{enumerate}

\chapter{Reference: row and column layout}
\label{ch:refman-rowcol}

\begin{aside}
Flex containers in two flavours. \code{h(...)} (or
\code{row}) lays children horizontally; \code{v(...)} (or
\code{column}) stacks them vertically.
\end{aside}

\section{Working example}

\begin{lstlisting}[language=C++]
#include <maya/maya.hpp>

using namespace maya;
using namespace maya::dsl;

struct RowCol {
    struct Model {};
    struct Quit {};
    using Msg = std::variant<Quit>;

    static Model init() { return {}; }
    static auto update(Model m, Msg)
        -> std::pair<Model, Cmd<Msg>> {
        return {m, Cmd<Msg>::quit()};
    }

    static Element view(const Model&) {
        return v(
            t<"a column of rows"> | Bold,
            blank_,
            h(
                t<"left">  | Fg<255, 200, 100>,
                spacer_,
                t<"right"> | Fg<100, 200, 255>
            ),
            h(
                t<"A">, text(" "),
                t<"B">, text(" "),
                t<"C">
            ),
            blank_,
            t<"q to quit"> | Dim
        ) | pad<1> | border_<Round>;
    }

    static auto subscribe(const Model&) -> Sub<Msg> {
        return key_map<Msg>({ {'q', Quit{}} });
    }
};

int main() { run<RowCol>({.title = "rowcol"}); }
\end{lstlisting}

\section{Notes}

\begin{itemize}[leftmargin=*]
  \item \code{spacer\_} expands to fill remaining space.
  \item \code{blank\_} is a blank line; equivalent to
    \code{text(" ")} in a column.
  \item You can nest \code{h(v(...))} and \code{v(h(...))}
    freely; layout is computed top-down.
  \item Both containers accept any number of children via
    parameter packs.
\end{itemize}

\section{Sizing}

\begin{itemize}[leftmargin=*]
  \item Default: children take their natural size.
  \item \code{flex<N>}: child gets \code{N} units of weight.
  \item \code{fixed\_width<N>}, \code{fixed\_height<N>}:
    hard size.
\end{itemize}

\section{Recap}

\begin{recap}
\begin{itemize}[leftmargin=*]
  \item \code{h} for rows, \code{v} for columns.
  \item \code{spacer\_} pushes; \code{blank\_} spaces.
  \item Nesting is free; sizing is via modifiers.
\end{itemize}
\end{recap}

\section*{Workshop}

\begin{enumerate}[leftmargin=*]
  \item Build a header bar with title left, date right.
  \item Make a three-column dashboard.
  \item Add a footer with hint text and a help shortcut.
\end{enumerate}

\chapter{Reference: padding, margin, border}
\label{ch:refman-padding}

\begin{aside}
Three decorations that change how an element sits in space.
Padding adds breathing room \emph{inside} a border; margin
pushes it away from siblings; border draws a line around it.
\end{aside}

\section{Forms}

\begin{itemize}[leftmargin=*]
  \item \code{pad<N>}: equal padding on all four sides.
  \item \code{pad<T, R, B, L>}: per-edge padding.
  \item \code{margin<N>}, \code{margin<T, R, B, L>}: same
    pattern.
  \item \code{border\_<Round>}, \code{border\_<Square>},
    \code{border\_<Heavy>}, \code{border\_<Double>}: line
    styles.
  \item \code{border\_<None>}: explicit no border (useful in
    conditional layouts).
\end{itemize}

\section{Working example}

\begin{lstlisting}[language=C++]
#include <maya/maya.hpp>

using namespace maya;
using namespace maya::dsl;

struct PadDemo {
    struct Model {}; struct Quit {};
    using Msg = std::variant<Quit>;

    static Model init() { return {}; }
    static auto update(Model m, Msg)
        -> std::pair<Model, Cmd<Msg>> {
        return {m, Cmd<Msg>::quit()};
    }

    static Element view(const Model&) {
        auto card = [](auto title, auto body) {
            return v(title | Bold, blank_, body)
                | pad<1> | border_<Round>;
        };

        return v(
            card(t<"round">,  t<"hello">),
            card(t<"square">, t<"hello">) | border_<Square>,
            card(t<"heavy">,  t<"hello">) | border_<Heavy>,
            card(t<"double">, t<"hello">) | border_<Double>,
            blank_,
            t<"q to quit"> | Dim
        );
    }

    static auto subscribe(const Model&) -> Sub<Msg> {
        return key_map<Msg>({ {'q', Quit{}} });
    }
};

int main() { run<PadDemo>({.title = "pad"}); }
\end{lstlisting}

\section{Notes}

\begin{itemize}[leftmargin=*]
  \item Apply with the pipe operator: \code{element | pad<2>}.
  \item Borders consume one cell on every side.
  \item Margin and padding stack; you can have both.
\end{itemize}

\section{Recap}

\begin{recap}
\begin{itemize}[leftmargin=*]
  \item \code{pad}, \code{margin}, \code{border\_}.
  \item Pipe to apply; stack for compound effects.
\end{itemize}
\end{recap}

\section*{Workshop}

\begin{enumerate}[leftmargin=*]
  \item Make a card with \code{pad<2, 4, 2, 4>}.
  \item Toggle the border style with a key press.
  \item Build a button: padded text inside a heavy border.
\end{enumerate}

\chapter{Reference: colour and style modifiers}
\label{ch:refman-style}

\begin{aside}
The full palette of inline style. Apply with the pipe
operator; chain as many as you like; the order is mostly
irrelevant.
\end{aside}

\section{Foreground / background}

\begin{itemize}[leftmargin=*]
  \item \code{Fg<R, G, B>}: 24-bit foreground.
  \item \code{Bg<R, G, B>}: 24-bit background.
  \item \code{Fg256<N>}, \code{Bg256<N>}: 256-colour
    palette index.
  \item \code{Fg16<Red>}, \code{Bg16<Blue>}: legacy 16-colour
    names.
\end{itemize}

\section{Weight and decoration}

\begin{itemize}[leftmargin=*]
  \item \code{Bold}, \code{Dim}, \code{Italic},
    \code{Underline}.
  \item \code{Blink}, \code{Reverse}, \code{Strike},
    \code{Hidden}.
\end{itemize}

\section{Working example}

\begin{lstlisting}[language=C++]
#include <maya/maya.hpp>

using namespace maya;
using namespace maya::dsl;

struct StyleDemo {
    struct Model {}; struct Quit {};
    using Msg = std::variant<Quit>;

    static Model init() { return {}; }
    static auto update(Model m, Msg)
        -> std::pair<Model, Cmd<Msg>> {
        return {m, Cmd<Msg>::quit()};
    }

    static Element view(const Model&) {
        return v(
            t<"bold">              | Bold,
            t<"dim">               | Dim,
            t<"italic">            | Italic,
            t<"underline">         | Underline,
            t<"strike">            | Strike,
            t<"reverse">           | Reverse,
            blank_,
            t<"red on black">      | Fg<255, 80, 80>,
            t<"green on dark">     | Fg<100, 220, 120>,
            t<"blue on dark">      | Fg<100, 180, 255>,
            t<"warning">           | Fg<20, 20, 20> | Bg<255, 220, 80>,
            blank_,
            t<"bold + italic + cyan"> | Bold | Italic | Fg<100, 220, 220>,
            blank_,
            t<"q to quit">         | Dim
        ) | pad<1> | border_<Round>;
    }

    static auto subscribe(const Model&) -> Sub<Msg> {
        return key_map<Msg>({ {'q', Quit{}} });
    }
};

int main() { run<StyleDemo>({.title = "style"}); }
\end{lstlisting}

\section{Notes}

\begin{itemize}[leftmargin=*]
  \item Modifiers compose left-to-right; later ones win on
    conflicts (e.g. two foregrounds).
  \item True colour is the default; the runtime
    automatically degrades to 256-colour where the terminal
    cannot do 24-bit.
  \item \code{Blink} is honoured rarely; treat as a hint.
\end{itemize}

\section{Recap}

\begin{recap}
\begin{itemize}[leftmargin=*]
  \item Fg / Bg in 24-bit, 256, or 16.
  \item Bold, Dim, Italic, Underline, Strike, Reverse.
  \item Chain with pipes.
\end{itemize}
\end{recap}

\section*{Workshop}

\begin{enumerate}[leftmargin=*]
  \item Build a status bar with three coloured pills.
  \item Use \code{Reverse} to highlight a selected row.
  \item Make a "danger" badge: bold, white-on-red.
\end{enumerate}

\chapter{Reference: alignment and centering}
\label{ch:refman-align}

\begin{aside}
Push content to a corner, the middle, or anywhere in between.
The alignment helpers wrap a child and shift it within the
slot they receive.
\end{aside}

\section{Forms}

\begin{itemize}[leftmargin=*]
  \item \code{center(child)}: middle on both axes.
  \item \code{align(x, y, child)}: \code{x, y} in [0, 1];
    \code{0} is left/top, \code{1} is right/bottom.
  \item \code{left(child)}, \code{right(child)},
    \code{top(child)}, \code{bottom(child)}: edge sugar.
\end{itemize}

\section{Working example}

\begin{lstlisting}[language=C++]
#include <maya/maya.hpp>

using namespace maya;
using namespace maya::dsl;

struct AlignDemo {
    struct Model {}; struct Quit {};
    using Msg = std::variant<Quit>;

    static Model init() { return {}; }
    static auto update(Model m, Msg)
        -> std::pair<Model, Cmd<Msg>> {
        return {m, Cmd<Msg>::quit()};
    }

    static Element view(const Model&) {
        return v(
            center(t<"centred"> | Bold) | fixed_height<5> | border_<Round>,
            blank_,
            h(
                left(t<"left">)   | flex<1> | border_<Round>,
                center(t<"mid">)  | flex<1> | border_<Round>,
                right(t<"right">) | flex<1> | border_<Round>
            ),
            blank_,
            t<"q to quit"> | Dim
        );
    }

    static auto subscribe(const Model&) -> Sub<Msg> {
        return key_map<Msg>({ {'q', Quit{}} });
    }
};

int main() { run<AlignDemo>({.title = "align"}); }
\end{lstlisting}

\section{Notes}

\begin{itemize}[leftmargin=*]
  \item Alignment is only visible if the wrapper has spare
    space. A naked \code{center} around a wider element is a
    no-op.
  \item \code{flex<1>} on siblings divides the row's spare
    width equally.
  \item Use \code{fixed\_height} to give alignment something
    to push against vertically.
\end{itemize}

\section{Recap}

\begin{recap}
\begin{itemize}[leftmargin=*]
  \item \code{center}, \code{align}, edge sugar.
  \item Visible only with extra room.
  \item Combine with \code{flex} for shared-width groups.
\end{itemize}
\end{recap}

\section*{Workshop}

\begin{enumerate}[leftmargin=*]
  \item Build a splash screen: centred title and subtitle.
  \item Align a status bar to the bottom.
  \item Use \code{align(0.5, 0.25, ...)} to put a banner at
    the top quarter.
\end{enumerate}

\chapter{Reference: subscriptions}
\label{ch:refman-sub}

\begin{aside}
Subscriptions are how the outside world reaches your update
function. Keys, timers, OS signals, file watches — every
recurring input is a \code{Sub<Msg>}.
\end{aside}

\section{Constructors}

\begin{itemize}[leftmargin=*]
  \item \code{Sub<Msg>::none()} — empty.
  \item \code{Sub<Msg>::batch(a, b, ...)} — combine several.
  \item \code{Sub<Msg>::every(duration, msg)} — periodic.
  \item \code{key\_map<Msg>(\{...\})} — declarative keys.
  \item \code{mouse\_map<Msg>(...)} — mouse events.
  \item \code{Sub<Msg>::on\_resize(fn)} — window resize.
  \item \code{Sub<Msg>::on\_signal(SIGUSR1, msg)} — OS signal.
\end{itemize}

\section{Working example}

\begin{lstlisting}[language=C++]
#include <maya/maya.hpp>
#include <chrono>

using namespace maya;
using namespace maya::dsl;

struct SubDemo {
    struct Model {
        int ticks = 0;
        int w = 0, h = 0;
    };
    struct Tick {};
    struct Resize { int w; int h; };
    struct Quit {};
    using Msg = std::variant<Tick, Resize, Quit>;

    static Model init() { return {}; }

    static auto update(Model m, Msg msg)
        -> std::pair<Model, Cmd<Msg>>
    {
        return std::visit(overload{
            [&](Tick) { m.ticks++; return std::pair{m, Cmd<Msg>{}}; },
            [&](Resize r) {
                m.w = r.w; m.h = r.h;
                return std::pair{m, Cmd<Msg>{}};
            },
            [](Quit) { return std::pair{Model{}, Cmd<Msg>::quit()}; },
        }, msg);
    }

    static Element view(const Model& m) {
        return v(
            text(std::format("ticks: {}", m.ticks)),
            text(std::format("size:  {}x{}", m.w, m.h)),
            blank_,
            t<"q to quit"> | Dim
        ) | pad<1> | border_<Round>;
    }

    static auto subscribe(const Model&) -> Sub<Msg> {
        return Sub<Msg>::batch(
            key_map<Msg>({ {'q', Quit{}} }),
            Sub<Msg>::every(std::chrono::milliseconds(500), Tick{}),
            Sub<Msg>::on_resize([](int w, int h) {
                return Resize{w, h};
            })
        );
    }
};

int main() { run<SubDemo>({.title = "sub"}); }
\end{lstlisting}

\section{Notes}

\begin{itemize}[leftmargin=*]
  \item \code{subscribe} is re-run every frame; you can
    return different subs based on Model state.
  \item Timers paused by returning a subscription \emph{without}
    that \code{every} entry.
  \item \code{batch} flattens nested batches.
\end{itemize}

\section{Recap}

\begin{recap}
\begin{itemize}[leftmargin=*]
  \item Keys, timers, resize, signals, mouse.
  \item Combine with \code{batch}.
  \item Conditionally include for pause/resume behaviour.
\end{itemize}
\end{recap}

\section*{Workshop}

\begin{enumerate}[leftmargin=*]
  \item Subscribe to \code{SIGINT} to clean up on Ctrl-C.
  \item Pause the timer when a flag in Model is false.
  \item Track the last key pressed in the Model.
\end{enumerate}

\chapter{Reference: \code{key\_map} keys}
\label{ch:refman-keymap}

\begin{aside}
Bind keys to messages declaratively. Both printable
characters and special keys go in the same map.
\end{aside}

\section{What you can put in the map}

\begin{itemize}[leftmargin=*]
  \item Single chars: \code{'q'}, \code{' '}, \code{'\textbackslash n'}.
  \item Special keys: \code{SpecialKey::Up},
    \code{Down}, \code{Left}, \code{Right},
    \code{Enter}, \code{Escape}, \code{Tab},
    \code{Backspace}, \code{Delete}, \code{Home},
    \code{End}, \code{PageUp}, \code{PageDown},
    \code{F1}..\code{F12}.
  \item Modifier combos: \code{ctrl('c')},
    \code{alt('x')}, \code{shift(SpecialKey::Up)}.
\end{itemize}

\section{Working example}

\begin{lstlisting}[language=C++]
#include <maya/maya.hpp>

using namespace maya;
using namespace maya::dsl;

struct KeyDemo {
    struct Model { std::string last; };
    struct Pressed { std::string label; };
    struct Quit {};
    using Msg = std::variant<Pressed, Quit>;

    static Model init() { return {}; }

    static auto update(Model m, Msg msg)
        -> std::pair<Model, Cmd<Msg>>
    {
        return std::visit(overload{
            [&](Pressed p) { m.last = p.label; return std::pair{m, Cmd<Msg>{}}; },
            [](Quit)       { return std::pair{Model{}, Cmd<Msg>::quit()}; },
        }, msg);
    }

    static Element view(const Model& m) {
        return v(
            t<"press a key">,
            blank_,
            text(std::format("last: {}", m.last)) | Bold,
            blank_,
            t<"q to quit"> | Dim
        ) | pad<1> | border_<Round>;
    }

    static auto subscribe(const Model&) -> Sub<Msg> {
        return key_map<Msg>({
            {'q',                  Quit{}},
            {' ',                  Pressed{"space"}},
            {SpecialKey::Up,       Pressed{"up"}},
            {SpecialKey::Down,     Pressed{"down"}},
            {SpecialKey::Enter,    Pressed{"enter"}},
            {SpecialKey::Escape,   Pressed{"escape"}},
            {SpecialKey::Tab,      Pressed{"tab"}},
            {ctrl('c'),            Quit{}},
            {ctrl('s'),            Pressed{"ctrl-s"}},
            {alt('x'),             Pressed{"alt-x"}},
            {SpecialKey::F1,       Pressed{"F1"}},
        });
    }
};

int main() { run<KeyDemo>({.title = "keys"}); }
\end{lstlisting}

\section{Notes}

\begin{itemize}[leftmargin=*]
  \item Order does not matter; lookup is by key value.
  \item Duplicate keys produce a runtime error at startup.
  \item Unhandled keys are silently dropped.
\end{itemize}

\section{Recap}

\begin{recap}
\begin{itemize}[leftmargin=*]
  \item Map literal keys, special keys, and modifier combos.
  \item One entry per key.
  \item Anything else is ignored.
\end{itemize}
\end{recap}

\section*{Workshop}

\begin{enumerate}[leftmargin=*]
  \item Bind \code{?} to show a help overlay.
  \item Bind \code{ctrl('r')} to reload.
  \item Bind every digit \code{0}--\code{9} to a tab switch.
\end{enumerate}

\chapter{Reference: timers and ticks}
\label{ch:refman-timers}

\begin{aside}
Periodic time-driven messages. Use them for animation,
polling, and anything else that needs to wake up on a
schedule.
\end{aside}

\section{Forms}

\begin{itemize}[leftmargin=*]
  \item \code{Sub<Msg>::every(duration, msg)}: send \code{msg}
    every \code{duration}.
  \item \code{Cmd<Msg>::after(duration, msg)}: send
    \code{msg} once, \code{duration} from now.
\end{itemize}

\section{Working example}

\begin{lstlisting}[language=C++]
#include <maya/maya.hpp>
#include <chrono>

using namespace maya;
using namespace maya::dsl;

struct Timer {
    struct Model {
        int seconds = 0;
        bool flashed = false;
    };
    struct Tick {};
    struct Flash {};
    struct ClearFlash {};
    struct Quit {};
    using Msg = std::variant<Tick, Flash, ClearFlash, Quit>;

    static Model init() { return {}; }

    static auto update(Model m, Msg msg)
        -> std::pair<Model, Cmd<Msg>>
    {
        return std::visit(overload{
            [&](Tick) {
                m.seconds++;
                return std::pair{m, Cmd<Msg>{}};
            },
            [&](Flash) {
                m.flashed = true;
                return std::pair{m,
                    Cmd<Msg>::after(std::chrono::milliseconds(500),
                                    ClearFlash{})};
            },
            [&](ClearFlash) {
                m.flashed = false;
                return std::pair{m, Cmd<Msg>{}};
            },
            [](Quit) {
                return std::pair{Model{}, Cmd<Msg>::quit()};
            },
        }, msg);
    }

    static Element view(const Model& m) {
        auto label = std::format("elapsed: {}s", m.seconds);
        return v(
            when(m.flashed,
                text(label) | Bold | Fg<255, 220, 80>,
                text(label) | Bold),
            blank_,
            t<"f flash, q quit"> | Dim
        ) | pad<1> | border_<Round>;
    }

    static auto subscribe(const Model&) -> Sub<Msg> {
        return Sub<Msg>::batch(
            key_map<Msg>({
                {'q', Quit{}},
                {'f', Flash{}},
            }),
            Sub<Msg>::every(std::chrono::seconds(1), Tick{})
        );
    }
};

int main() { run<Timer>({.title = "timer"}); }
\end{lstlisting}

\section{Notes}

\begin{itemize}[leftmargin=*]
  \item Resolution is bounded by your \code{fps}; finer
    timers fire on the next available frame.
  \item \code{after} is "fire once" — don't use it for
    repeats.
  \item Timers do not drift; they're scheduled against a
    monotonic clock.
\end{itemize}

\section{Recap}

\begin{recap}
\begin{itemize}[leftmargin=*]
  \item \code{every} for repeats, \code{after} for one-shots.
  \item Resolution is frame-bounded.
  \item Use \code{after} to clear transient highlights.
\end{itemize}
\end{recap}

\section*{Workshop}

\begin{enumerate}[leftmargin=*]
  \item Animate a spinner using \code{every(80ms, ...)}.
  \item Show a "saved" toast for two seconds, then clear.
  \item Build a Pomodoro timer with 25-minute work blocks.
\end{enumerate}

\chapter{Reference: commands}
\label{ch:refman-cmd}

\begin{aside}
A \code{Cmd<Msg>} is a description of an effect to perform.
It is returned from \code{update}; the runtime performs the
effect and, if it produces a result, dispatches a message.
\end{aside}

\section{Constructors}

\begin{itemize}[leftmargin=*]
  \item \code{Cmd<Msg>\{\}} — no effect.
  \item \code{Cmd<Msg>::none()} — same; explicit form.
  \item \code{Cmd<Msg>::quit()} — exit the program.
  \item \code{Cmd<Msg>::after(d, msg)} — delayed message.
  \item \code{Cmd<Msg>::batch(a, b, ...)} — combine.
  \item \code{Cmd<Msg>::task(callable, then\_msg)} — run on a
    worker, dispatch result.
  \item \code{Cmd<Msg>::http\_get(url, then\_msg)} — async
    HTTP.
  \item \code{Cmd<Msg>::read\_file(path, then\_msg)} — async
    file read.
\end{itemize}

\section{Working example}

\begin{lstlisting}[language=C++]
#include <maya/maya.hpp>
#include <chrono>

using namespace maya;
using namespace maya::dsl;

struct CmdDemo {
    struct Model { int n = 0; std::string status; };
    struct Bump {};
    struct DelayedBump {};
    struct Quit {};
    using Msg = std::variant<Bump, DelayedBump, Quit>;

    static Model init() { return {}; }

    static auto update(Model m, Msg msg)
        -> std::pair<Model, Cmd<Msg>>
    {
        return std::visit(overload{
            [&](Bump) {
                m.n++;
                m.status = "bumped";
                return std::pair{m, Cmd<Msg>{}};
            },
            [&](DelayedBump) {
                m.status = "scheduled...";
                return std::pair{m,
                    Cmd<Msg>::after(std::chrono::seconds(1),
                                    Bump{})};
            },
            [](Quit) {
                return std::pair{Model{}, Cmd<Msg>::quit()};
            },
        }, msg);
    }

    static Element view(const Model& m) {
        return v(
            text(std::format("n = {}", m.n)) | Bold,
            text(m.status) | Dim,
            blank_,
            t<"b bump, d delayed, q quit"> | Dim
        ) | pad<1> | border_<Round>;
    }

    static auto subscribe(const Model&) -> Sub<Msg> {
        return key_map<Msg>({
            {'q', Quit{}},
            {'b', Bump{}},
            {'d', DelayedBump{}},
        });
    }
};

int main() { run<CmdDemo>({.title = "cmd"}); }
\end{lstlisting}

\section{Notes}

\begin{itemize}[leftmargin=*]
  \item Commands are values. You build them, return them,
    never invoke them yourself.
  \item Order inside \code{batch} is preserved.
  \item \code{quit} is irrevocable: returning it once tears
    the loop down.
\end{itemize}

\section{Recap}

\begin{recap}
\begin{itemize}[leftmargin=*]
  \item Effects are data; the runtime performs them.
  \item Common constructors: none, quit, after, batch, task,
    http\_get, read\_file.
\end{itemize}
\end{recap}

\section*{Workshop}

\begin{enumerate}[leftmargin=*]
  \item Implement a button that triggers two commands at
    once with \code{batch}.
  \item Use \code{task} to compute a fibonacci on a worker
    thread; show the result.
  \item Trigger a \code{quit} two seconds after a key press.
\end{enumerate}

\chapter{Reference: HTTP and tasks}
\label{ch:refman-task}

\begin{aside}
Anything that blocks belongs in a command. Files, network,
process spawns — return a \code{task} or \code{http\_get}
and \maya{} will run it off the loop, then post the result
back as a message.
\end{aside}

\section{Forms}

\begin{itemize}[leftmargin=*]
  \item \code{Cmd<Msg>::task(callable, then\_msg\_fn)}: run
    \code{callable()} on a worker, pass the result to
    \code{then\_msg\_fn} to make a \code{Msg}.
  \item \code{Cmd<Msg>::http\_get(url, then\_msg\_fn)}: GET
    request, dispatch a message with the body or error.
  \item \code{Cmd<Msg>::read\_file(path, then\_msg\_fn)}:
    async file read.
\end{itemize}

\section{Working example}

\begin{lstlisting}[language=C++]
#include <maya/maya.hpp>
#include <thread>

using namespace maya;
using namespace maya::dsl;

struct TaskDemo {
    struct Model {
        std::string status = "idle";
        long answer = 0;
    };
    struct Start {};
    struct Done { long value; };
    struct Quit {};
    using Msg = std::variant<Start, Done, Quit>;

    static Model init() { return {}; }

    static auto update(Model m, Msg msg)
        -> std::pair<Model, Cmd<Msg>>
    {
        return std::visit(overload{
            [&](Start) {
                m.status = "computing...";
                auto compute = []() -> long {
                    std::this_thread::sleep_for(
                        std::chrono::seconds(2));
                    long sum = 0;
                    for (long i = 1; i <= 1000000; ++i) sum += i;
                    return sum;
                };
                return std::pair{m,
                    Cmd<Msg>::task(compute,
                        [](long v) { return Done{v}; })};
            },
            [&](Done d) {
                m.status = "done";
                m.answer = d.value;
                return std::pair{m, Cmd<Msg>{}};
            },
            [](Quit) {
                return std::pair{Model{}, Cmd<Msg>::quit()};
            },
        }, msg);
    }

    static Element view(const Model& m) {
        return v(
            text(m.status) | Bold,
            text(std::format("answer: {}", m.answer)),
            blank_,
            t<"s start, q quit"> | Dim
        ) | pad<1> | border_<Round>;
    }

    static auto subscribe(const Model&) -> Sub<Msg> {
        return key_map<Msg>({
            {'q', Quit{}},
            {'s', Start{}},
        });
    }
};

int main() { run<TaskDemo>({.title = "task"}); }
\end{lstlisting}

\section{Notes}

\begin{itemize}[leftmargin=*]
  \item The worker pool is shared; long tasks block other
    tasks, not the UI.
  \item Exceptions in the callable become a \code{Done}
    with a default-constructed value plus an error flag if
    your \code{Msg} carries one.
  \item Never touch the Model from inside the task. Only
    the returned message may mutate state.
\end{itemize}

\section{Recap}

\begin{recap}
\begin{itemize}[leftmargin=*]
  \item \code{task}, \code{http\_get}, \code{read\_file}.
  \item Blocking work goes off-loop; result returns as a
    message.
  \item No shared state with the worker.
\end{itemize}
\end{recap}

\section*{Workshop}

\begin{enumerate}[leftmargin=*]
  \item Spawn five tasks at once; show progress as each
    finishes.
  \item Fetch a URL and render the first 200 bytes.
  \item Read a file and count its lines off-loop.
\end{enumerate}

\chapter{Reference: mouse}
\label{ch:refman-mouse}

\begin{aside}
Mouse support is opt-in. Set \code{RunConfig::mouse = true},
add a mouse handler in \code{subscribe}, and clicks /
scrolls / drags arrive as messages.
\end{aside}

\section{Event types}

\begin{itemize}[leftmargin=*]
  \item \code{MouseButton::Left}, \code{Middle},
    \code{Right}.
  \item \code{MouseKind::Press}, \code{Release},
    \code{Drag}, \code{Move}.
  \item Wheel: \code{MouseButton::WheelUp},
    \code{WheelDown}.
\end{itemize}

\section{Working example}

\begin{lstlisting}[language=C++]
#include <maya/maya.hpp>

using namespace maya;
using namespace maya::dsl;

struct MouseDemo {
    struct Model {
        int x = 0, y = 0;
        std::string last = "(none)";
    };
    struct MouseEvt { int x; int y; std::string label; };
    struct Quit {};
    using Msg = std::variant<MouseEvt, Quit>;

    static Model init() { return {}; }

    static auto update(Model m, Msg msg)
        -> std::pair<Model, Cmd<Msg>>
    {
        return std::visit(overload{
            [&](MouseEvt e) {
                m.x = e.x; m.y = e.y; m.last = e.label;
                return std::pair{m, Cmd<Msg>{}};
            },
            [](Quit) {
                return std::pair{Model{}, Cmd<Msg>::quit()};
            },
        }, msg);
    }

    static Element view(const Model& m) {
        return v(
            text(std::format("pos: ({}, {})", m.x, m.y)),
            text(std::format("evt: {}", m.last)),
            blank_,
            t<"click anywhere, q to quit"> | Dim
        ) | pad<1> | border_<Round>;
    }

    static auto subscribe(const Model&) -> Sub<Msg> {
        return Sub<Msg>::batch(
            key_map<Msg>({ {'q', Quit{}} }),
            mouse_map<Msg>([](MouseEvent e) -> Msg {
                std::string lbl;
                switch (e.kind) {
                    case MouseKind::Press:   lbl = "press"; break;
                    case MouseKind::Release: lbl = "release"; break;
                    case MouseKind::Drag:    lbl = "drag"; break;
                    case MouseKind::Move:    lbl = "move"; break;
                }
                return MouseEvt{e.x, e.y, lbl};
            })
        );
    }
};

int main() {
    return run<MouseDemo>({
        .title = "mouse",
        .mouse = true,
    });
}
\end{lstlisting}

\section{Notes}

\begin{itemize}[leftmargin=*]
  \item Coordinates are zero-based; \code{(0, 0)} is the
    top-left cell.
  \item Some terminals require additional config to forward
    mouse events through tmux/screen.
  \item Selection-by-drag in the terminal may conflict with
    \maya{} mouse capture; users hold Shift to override.
\end{itemize}

\section{Recap}

\begin{recap}
\begin{itemize}[leftmargin=*]
  \item Opt in via \code{RunConfig}.
  \item Press, release, drag, move, wheel.
  \item Coordinates in terminal cells.
\end{itemize}
\end{recap}

\section*{Workshop}

\begin{enumerate}[leftmargin=*]
  \item Build a clickable button: highlight on press,
    fire a command on release.
  \item Implement drag-to-scroll on a long list.
  \item Show a tooltip on hover.
\end{enumerate}

\chapter{Reference: \code{when} and conditionals}
\label{ch:refman-when}

\begin{aside}
A tiny helper for branching inside an element tree without
breaking the expression chain.
\end{aside}

\section{Forms}

\begin{itemize}[leftmargin=*]
  \item \code{when(cond, then\_element)} — render
    \code{then\_element} if \code{cond}, otherwise nothing.
  \item \code{when(cond, then, else)} — ternary form.
\end{itemize}

\section{Working example}

\begin{lstlisting}[language=C++]
#include <maya/maya.hpp>

using namespace maya;
using namespace maya::dsl;

struct WhenDemo {
    struct Model { bool on = false; int count = 0; };
    struct Toggle {};
    struct Bump {};
    struct Quit {};
    using Msg = std::variant<Toggle, Bump, Quit>;

    static Model init() { return {}; }

    static auto update(Model m, Msg msg)
        -> std::pair<Model, Cmd<Msg>>
    {
        return std::visit(overload{
            [&](Toggle) { m.on = !m.on; return std::pair{m, Cmd<Msg>{}}; },
            [&](Bump)   { m.count++;    return std::pair{m, Cmd<Msg>{}}; },
            [](Quit)    { return std::pair{Model{}, Cmd<Msg>::quit()}; },
        }, msg);
    }

    static Element view(const Model& m) {
        return v(
            when(m.on,
                t<"ON">  | Bold | Fg<100, 255, 100>,
                t<"OFF"> | Bold | Fg<255, 100, 100>),
            blank_,
            when(m.count > 0,
                text(std::format("clicks: {}", m.count))),
            when(m.count >= 5,
                t<"high five!"> | Bold | Fg<255, 220, 80>),
            blank_,
            t<"t toggle, b bump, q quit"> | Dim
        ) | pad<1> | border_<Round>;
    }

    static auto subscribe(const Model&) -> Sub<Msg> {
        return key_map<Msg>({
            {'q', Quit{}},
            {'t', Toggle{}},
            {'b', Bump{}},
        });
    }
};

int main() { run<WhenDemo>({.title = "when"}); }
\end{lstlisting}

\section{Notes}

\begin{itemize}[leftmargin=*]
  \item \code{when(false, e)} contributes nothing to layout.
  \item Both branches are \emph{constructed} either way;
    pick the cheap form if construction is expensive.
  \item For long branches, use a helper function and inline
    the call.
\end{itemize}

\section{Recap}

\begin{recap}
\begin{itemize}[leftmargin=*]
  \item \code{when} for inline branches.
  \item Two-arg "show if", three-arg "ternary".
  \item Both branches construct; only one renders.
\end{itemize}
\end{recap}

\section*{Workshop}

\begin{enumerate}[leftmargin=*]
  \item Use \code{when} to hide a footer when the window is
    too short.
  \item Render different views per "tab" using a sequence of
    \code{when}s.
  \item Replace a long \code{if/else} chain with nested
    \code{when}s.
\end{enumerate}

\chapter{Reference: lists and selection}
\label{ch:refman-list}

\begin{aside}
A scrolling list of items with a highlighted selection.
Build it from rows and a model index; \maya{} ships nothing
fancier than that because nothing fancier is needed.
\end{aside}

\section{Working example}

\begin{lstlisting}[language=C++]
#include <maya/maya.hpp>
#include <vector>

using namespace maya;
using namespace maya::dsl;

struct ListDemo {
    struct Model {
        std::vector<std::string> items{
            "alpha", "beta", "gamma", "delta",
            "epsilon", "zeta", "eta", "theta"
        };
        int selected = 0;
    };
    struct Up {};
    struct Down {};
    struct Choose {};
    struct Quit {};
    using Msg = std::variant<Up, Down, Choose, Quit>;

    static Model init() { return {}; }

    static auto update(Model m, Msg msg)
        -> std::pair<Model, Cmd<Msg>>
    {
        return std::visit(overload{
            [&](Up) {
                m.selected = std::max(0, m.selected - 1);
                return std::pair{m, Cmd<Msg>{}};
            },
            [&](Down) {
                int n = static_cast<int>(m.items.size());
                m.selected = std::min(n - 1, m.selected + 1);
                return std::pair{m, Cmd<Msg>{}};
            },
            [&](Choose) { return std::pair{m, Cmd<Msg>{}}; },
            [](Quit) {
                return std::pair{Model{}, Cmd<Msg>::quit()};
            },
        }, msg);
    }

    static Element view(const Model& m) {
        std::vector<Element> rows;
        for (size_t i = 0; i < m.items.size(); ++i) {
            auto label = text(m.items[i]);
            if (static_cast<int>(i) == m.selected) {
                rows.push_back(
                    h(t<"> ">, label | Bold) | Bg<40, 80, 140>);
            } else {
                rows.push_back(h(t<"  ">, label));
            }
        }
        return v(
            t<"pick one"> | Bold,
            blank_,
            v(rows),
            blank_,
            t<"up/down move, enter choose, q quit"> | Dim
        ) | pad<1> | border_<Round>;
    }

    static auto subscribe(const Model&) -> Sub<Msg> {
        return key_map<Msg>({
            {'q', Quit{}},
            {SpecialKey::Up,    Up{}},
            {SpecialKey::Down,  Down{}},
            {SpecialKey::Enter, Choose{}},
        });
    }
};

int main() { run<ListDemo>({.title = "list"}); }
\end{lstlisting}

\section{Notes}

\begin{itemize}[leftmargin=*]
  \item For long lists, clip a window of items around the
    selection; rendering off-screen rows costs.
  \item Use \code{Bg<...>} on the selected row; readers
    parse highlight faster than markers.
  \item Pair with \code{home}/\code{end} for jump-to-edges.
\end{itemize}

\section{Recap}

\begin{recap}
\begin{itemize}[leftmargin=*]
  \item Build rows; mark one selected.
  \item Up / Down adjust a single index.
  \item Clip to a window for large lists.
\end{itemize}
\end{recap}

\section*{Workshop}

\begin{enumerate}[leftmargin=*]
  \item Add a filter input that prunes the visible items.
  \item Show the index and total in the footer.
  \item Implement multi-select with \code{space} to toggle.
\end{enumerate}

\chapter{Reference: text input}
\label{ch:refman-input}

\begin{aside}
A single-line editable buffer. Maya does not bundle a fully
featured editor widget; the example below is the canonical
60-line text input that covers the 95\% case.
\end{aside}

\section{Working example}

\begin{lstlisting}[language=C++]
#include <maya/maya.hpp>
#include <string>

using namespace maya;
using namespace maya::dsl;

struct Input {
    struct Model { std::string buffer; };
    struct Type { char c; };
    struct Back {};
    struct Submit {};
    struct Quit {};
    using Msg = std::variant<Type, Back, Submit, Quit>;

    static Model init() { return {}; }

    static auto update(Model m, Msg msg)
        -> std::pair<Model, Cmd<Msg>>
    {
        return std::visit(overload{
            [&](Type t) {
                m.buffer.push_back(t.c);
                return std::pair{m, Cmd<Msg>{}};
            },
            [&](Back) {
                if (!m.buffer.empty()) m.buffer.pop_back();
                return std::pair{m, Cmd<Msg>{}};
            },
            [&](Submit) { return std::pair{m, Cmd<Msg>{}}; },
            [](Quit) {
                return std::pair{Model{}, Cmd<Msg>::quit()};
            },
        }, msg);
    }

    static Element view(const Model& m) {
        return v(
            t<"name:"> | Bold,
            blank_,
            h(text(m.buffer), text("_") | Blink),
            blank_,
            t<"enter submit, esc quit"> | Dim
        ) | pad<1> | border_<Round> | fixed_width<40>;
    }

    static auto subscribe(const Model&) -> Sub<Msg> {
        // Build a key map covering every printable char.
        std::vector<std::pair<KeyEvent, Msg>> binds;
        for (char c = ' '; c < 127; ++c) {
            binds.push_back({c, Type{c}});
        }
        binds.push_back({SpecialKey::Backspace, Back{}});
        binds.push_back({SpecialKey::Enter,     Submit{}});
        binds.push_back({SpecialKey::Escape,    Quit{}});
        return key_map<Msg>(std::move(binds));
    }
};

int main() { run<Input>({.title = "input"}); }
\end{lstlisting}

\section{Notes}

\begin{itemize}[leftmargin=*]
  \item The cursor here is just a blinking underscore in the
    text; for a real cursor, position one with absolute
    output.
  \item Real apps separate input model from the rest of the
    Model so resets are cheap.
  \item For multi-line input, store a \code{std::vector<std::string>}.
\end{itemize}

\section{Recap}

\begin{recap}
\begin{itemize}[leftmargin=*]
  \item Roll your own input; it's small.
  \item One char per \code{Type} message; one byte per
    \code{Back}.
  \item Submit on Enter; cancel on Escape.
\end{itemize}
\end{recap}

\section*{Workshop}

\begin{enumerate}[leftmargin=*]
  \item Add left/right arrows to move a real caret.
  \item Implement word-wise backspace (\code{ctrl-w}).
  \item Add a placeholder shown when the buffer is empty.
\end{enumerate}

\chapter{Reference: progress bars}
\label{ch:refman-progress}

\begin{aside}
A determinate progress bar is a row of filled and empty
cells. Build it once; reuse forever.
\end{aside}

\section{Working example}

\begin{lstlisting}[language=C++]
#include <maya/maya.hpp>
#include <chrono>

using namespace maya;
using namespace maya::dsl;

namespace {
Element progress(float frac, int width) {
    frac = std::clamp(frac, 0.f, 1.f);
    int filled = static_cast<int>(frac * width + 0.5f);
    std::string bar;
    for (int i = 0; i < width; ++i) {
        bar += (i < filled) ? "#" : ".";
    }
    return text(bar);
}
}

struct ProgDemo {
    struct Model { float frac = 0.f; };
    struct Tick {};
    struct Quit {};
    using Msg = std::variant<Tick, Quit>;

    static Model init() { return {}; }

    static auto update(Model m, Msg msg)
        -> std::pair<Model, Cmd<Msg>>
    {
        return std::visit(overload{
            [&](Tick) {
                m.frac = std::min(1.f, m.frac + 0.01f);
                return std::pair{m, Cmd<Msg>{}};
            },
            [](Quit) {
                return std::pair{Model{}, Cmd<Msg>::quit()};
            },
        }, msg);
    }

    static Element view(const Model& m) {
        auto pct = std::format("{:>3}\%",
                               static_cast<int>(m.frac * 100));
        return v(
            t<"downloading"> | Bold,
            blank_,
            h(progress(m.frac, 30) | Fg<100, 220, 160>,
              text(" "), text(pct) | Bold),
            blank_,
            t<"q to quit"> | Dim
        ) | pad<1> | border_<Round>;
    }

    static auto subscribe(const Model& m) -> Sub<Msg> {
        auto keys = key_map<Msg>({ {'q', Quit{}} });
        if (m.frac < 1.f) {
            return Sub<Msg>::batch(std::move(keys),
                Sub<Msg>::every(std::chrono::milliseconds(50),
                                Tick{}));
        }
        return keys;
    }
};

int main() { run<ProgDemo>({.title = "progress"}); }
\end{lstlisting}

\section{Notes}

\begin{itemize}[leftmargin=*]
  \item Use unicode block characters (\code{U+2588} and
    friends) for higher-resolution progress; the eighths
    blocks give 8x the granularity at the same width.
  \item For indeterminate progress, animate a moving
    highlight inside a fixed-width row.
  \item Stop the tick subscription when work completes.
\end{itemize}

\section{Recap}

\begin{recap}
\begin{itemize}[leftmargin=*]
  \item Build a bar from filled and empty cells.
  \item Compute fill from a clamped fraction.
  \item Stop ticking when done.
\end{itemize}
\end{recap}

\section*{Workshop}

\begin{enumerate}[leftmargin=*]
  \item Show ETA next to the bar.
  \item Use eighths blocks for sub-cell precision.
  \item Build a stacked bar (multiple segments).
\end{enumerate}

\chapter{Reference: spinners}
\label{ch:refman-spinner}

\begin{aside}
A spinner is one character cycled at a fixed rate. The whole
widget is a frame counter and a string of glyphs.
\end{aside}

\section{Working example}

\begin{lstlisting}[language=C++]
#include <maya/maya.hpp>
#include <chrono>
#include <string_view>

using namespace maya;
using namespace maya::dsl;

namespace {
constexpr std::string_view frames[] = {
    "|", "/", "-", "\\\\",
};
}

struct Spin {
    struct Model { int frame = 0; bool busy = true; };
    struct Tick {};
    struct Toggle {};
    struct Quit {};
    using Msg = std::variant<Tick, Toggle, Quit>;

    static Model init() { return {}; }

    static auto update(Model m, Msg msg)
        -> std::pair<Model, Cmd<Msg>>
    {
        return std::visit(overload{
            [&](Tick) {
                m.frame = (m.frame + 1) % 4;
                return std::pair{m, Cmd<Msg>{}};
            },
            [&](Toggle) {
                m.busy = !m.busy;
                return std::pair{m, Cmd<Msg>{}};
            },
            [](Quit) {
                return std::pair{Model{}, Cmd<Msg>::quit()};
            },
        }, msg);
    }

    static Element view(const Model& m) {
        auto glyph = std::string(frames[m.frame]);
        return v(
            h(
                when(m.busy,
                    text(glyph) | Fg<100, 200, 255>,
                    t<"*">      | Fg<100, 220, 120>),
                text(" "),
                text(m.busy ? "working" : "idle") | Bold
            ),
            blank_,
            t<"space toggle, q quit"> | Dim
        ) | pad<1> | border_<Round>;
    }

    static auto subscribe(const Model& m) -> Sub<Msg> {
        auto keys = key_map<Msg>({
            {'q', Quit{}},
            {' ', Toggle{}},
        });
        if (m.busy) {
            return Sub<Msg>::batch(std::move(keys),
                Sub<Msg>::every(std::chrono::milliseconds(80),
                                Tick{}));
        }
        return keys;
    }
};

int main() { run<Spin>({.title = "spinner"}); }
\end{lstlisting}

\section{Frame sets}

\begin{itemize}[leftmargin=*]
  \item ASCII: \code{|}, \code{/}, \code{-}, \code{\textbackslash}.
  \item Braille: \code{U+2800}--\code{U+28FF}, eight-frame
    cycles look smooth.
  \item Dots: \code{.}, \code{..}, \code{...} for a slower
    pulse.
\end{itemize}

\section{Recap}

\begin{recap}
\begin{itemize}[leftmargin=*]
  \item Frame counter modulo frame count.
  \item Tick at 60--120 ms for human-pleasant motion.
  \item Stop ticking when work finishes.
\end{itemize}
\end{recap}

\section*{Workshop}

\begin{enumerate}[leftmargin=*]
  \item Swap the ASCII frames for braille dots.
  \item Render three spinners with different periods.
  \item Show a green tick instead of the spinner once work
    is done.
\end{enumerate}

\chapter{Reference: tables}
\label{ch:refman-table}

\begin{aside}
A table is a column of rows of columns. The expressive power
is in the helpers: column widths, header styling, zebra
stripes, alignment.
\end{aside}

\section{Working example}

\begin{lstlisting}[language=C++]
#include <maya/maya.hpp>
#include <vector>

using namespace maya;
using namespace maya::dsl;

struct TableDemo {
    struct Row { std::string name; int age; std::string job; };
    struct Model {
        std::vector<Row> rows{
            {"ada",    36, "mathematician"},
            {"alan",   42, "logician"},
            {"grace",  85, "compiler"},
            {"linus",  54, "kernel"},
        };
    };
    struct Quit {};
    using Msg = std::variant<Quit>;

    static Model init() { return {}; }
    static auto update(Model m, Msg)
        -> std::pair<Model, Cmd<Msg>> {
        return {m, Cmd<Msg>::quit()};
    }

    static Element view(const Model& m) {
        auto cell = [](std::string s, int w) {
            s.resize(w, ' ');
            return text(s);
        };

        auto header = h(
            cell("name", 10) | Bold,
            cell("age",   5) | Bold,
            cell("job",  16) | Bold
        );

        std::vector<Element> body;
        bool zebra = false;
        for (auto& r : m.rows) {
            auto row = h(
                cell(r.name,         10),
                cell(std::to_string(r.age), 5),
                cell(r.job,          16)
            );
            if (zebra) row = row | Bg<30, 30, 36>;
            zebra = !zebra;
            body.push_back(row);
        }

        return v(
            header,
            v(body),
            blank_,
            t<"q to quit"> | Dim
        ) | pad<1> | border_<Round>;
    }

    static auto subscribe(const Model&) -> Sub<Msg> {
        return key_map<Msg>({ {'q', Quit{}} });
    }
};

int main() { run<TableDemo>({.title = "table"}); }
\end{lstlisting}

\section{Notes}

\begin{itemize}[leftmargin=*]
  \item Pad cells to a fixed width before passing to
    \code{text}; uniform widths make rows align.
  \item Use \code{Bg<...>} for zebra stripes; readers track
    rows by the stripe.
  \item For sortable tables, sort the underlying vector;
    don't try to sort in render.
\end{itemize}

\section{Recap}

\begin{recap}
\begin{itemize}[leftmargin=*]
  \item Build cells, pad to fixed width.
  \item Header row above, body below.
  \item Zebra stripes with alternating background.
\end{itemize}
\end{recap}

\section*{Workshop}

\begin{enumerate}[leftmargin=*]
  \item Add a sort key; press a column letter to sort.
  \item Highlight the row under the cursor.
  \item Add a footer with the row count.
\end{enumerate}

\chapter{Reference: tabs}
\label{ch:refman-tabs}

\begin{aside}
Tabs partition a single content area into named panels.
Store the active index in the Model; render the matching
content; show the labels above.
\end{aside}

\section{Working example}

\begin{lstlisting}[language=C++]
#include <maya/maya.hpp>
#include <vector>

using namespace maya;
using namespace maya::dsl;

struct TabsDemo {
    struct Model {
        std::vector<std::string> labels{"info", "logs", "about"};
        int active = 0;
    };
    struct Next {};
    struct Prev {};
    struct Pick { int i; };
    struct Quit {};
    using Msg = std::variant<Next, Prev, Pick, Quit>;

    static Model init() { return {}; }

    static auto update(Model m, Msg msg)
        -> std::pair<Model, Cmd<Msg>>
    {
        int n = static_cast<int>(m.labels.size());
        return std::visit(overload{
            [&](Next) { m.active = (m.active + 1) % n; return std::pair{m, Cmd<Msg>{}}; },
            [&](Prev) { m.active = (m.active - 1 + n) % n; return std::pair{m, Cmd<Msg>{}}; },
            [&](Pick p) {
                if (p.i >= 0 && p.i < n) m.active = p.i;
                return std::pair{m, Cmd<Msg>{}};
            },
            [](Quit) {
                return std::pair{Model{}, Cmd<Msg>::quit()};
            },
        }, msg);
    }

    static Element view(const Model& m) {
        std::vector<Element> tabs;
        for (size_t i = 0; i < m.labels.size(); ++i) {
            auto label = text(" " + m.labels[i] + " ");
            if (static_cast<int>(i) == m.active) {
                tabs.push_back(label | Bold | Bg<60, 100, 180>);
            } else {
                tabs.push_back(label | Dim);
            }
        }

        Element body;
        switch (m.active) {
            case 0: body = t<"info: hello, world.">; break;
            case 1: body = t<"logs: nothing to report.">; break;
            default: body = t<"about: built with maya.">; break;
        }

        return v(
            h(tabs),
            blank_,
            body,
            blank_,
            t<"tab/shift-tab cycle, 1-3 pick, q quit"> | Dim
        ) | pad<1> | border_<Round>;
    }

    static auto subscribe(const Model&) -> Sub<Msg> {
        return key_map<Msg>({
            {'q',                Quit{}},
            {SpecialKey::Tab,    Next{}},
            {shift(SpecialKey::Tab), Prev{}},
            {'1', Pick{0}},
            {'2', Pick{1}},
            {'3', Pick{2}},
        });
    }
};

int main() { run<TabsDemo>({.title = "tabs"}); }
\end{lstlisting}

\section{Notes}

\begin{itemize}[leftmargin=*]
  \item Use a clear visual distinction between active and
    inactive tabs (background colour beats underline).
  \item Always include keyboard shortcuts; mouse-only tabs
    fail half your audience.
  \item For many tabs, scroll the label row horizontally.
\end{itemize}

\section{Recap}

\begin{recap}
\begin{itemize}[leftmargin=*]
  \item One active index; one switch in render.
  \item Tab to cycle; digits to jump.
  \item Bold + background highlights the active tab.
\end{itemize}
\end{recap}

\section*{Workshop}

\begin{enumerate}[leftmargin=*]
  \item Wrap each tab body in a separate function.
  \item Add a close-tab key for dynamic tab sets.
  \item Persist the active tab across runs.
\end{enumerate}

\chapter{Reference: modal dialogs}
\label{ch:refman-modal}

\begin{aside}
A modal is a centred panel rendered on top of the background.
Implement it as a flag in the Model and a top-level branch in
\code{view}.
\end{aside}

\section{Working example}

\begin{lstlisting}[language=C++]
#include <maya/maya.hpp>

using namespace maya;
using namespace maya::dsl;

struct Modal {
    struct Model { bool open = false; int count = 0; };
    struct Open {};
    struct Close {};
    struct Confirm {};
    struct Quit {};
    using Msg = std::variant<Open, Close, Confirm, Quit>;

    static Model init() { return {}; }

    static auto update(Model m, Msg msg)
        -> std::pair<Model, Cmd<Msg>>
    {
        return std::visit(overload{
            [&](Open)    { m.open = true;  return std::pair{m, Cmd<Msg>{}}; },
            [&](Close)   { m.open = false; return std::pair{m, Cmd<Msg>{}}; },
            [&](Confirm) {
                m.count++;
                m.open = false;
                return std::pair{m, Cmd<Msg>{}};
            },
            [](Quit) {
                return std::pair{Model{}, Cmd<Msg>::quit()};
            },
        }, msg);
    }

    static Element view(const Model& m) {
        auto base = v(
            text(std::format("confirmed: {}", m.count)) | Bold,
            blank_,
            t<"o open dialog, q quit"> | Dim
        ) | pad<1> | border_<Round>;

        if (!m.open) return base;

        auto dialog = v(
            t<"are you sure?"> | Bold,
            blank_,
            t<"y confirm, n cancel"> | Dim
        ) | pad<2> | border_<Double> | Bg<20, 20, 28>;

        return v(base, blank_, center(dialog));
    }

    static auto subscribe(const Model& m) -> Sub<Msg> {
        if (m.open) {
            return key_map<Msg>({
                {'y', Confirm{}},
                {'n', Close{}},
                {SpecialKey::Escape, Close{}},
            });
        }
        return key_map<Msg>({
            {'q', Quit{}},
            {'o', Open{}},
        });
    }
};

int main() { run<Modal>({.title = "modal"}); }
\end{lstlisting}

\section{Notes}

\begin{itemize}[leftmargin=*]
  \item Subscriptions branch on \code{open} so the dialog
    has its own key map; the background can't fight it.
  \item Background dim is optional but appreciated; render
    the base view with reduced contrast when open.
  \item Always provide an Escape route.
\end{itemize}

\section{Recap}

\begin{recap}
\begin{itemize}[leftmargin=*]
  \item One flag in the Model.
  \item Branch \code{view} and \code{subscribe} on it.
  \item Escape, always.
\end{itemize}
\end{recap}

\section*{Workshop}

\begin{enumerate}[leftmargin=*]
  \item Add a text input inside the dialog.
  \item Stack a second dialog ("are you really sure?").
  \item Dim the background when the dialog is open.
\end{enumerate}

\chapter{Reference: scroll regions}
\label{ch:refman-scroll}

\begin{aside}
Long content in a small window means a scrollable region.
Keep an offset in the Model; slice the content by the
offset; render only the visible window.
\end{aside}

\section{Working example}

\begin{lstlisting}[language=C++]
#include <maya/maya.hpp>
#include <vector>

using namespace maya;
using namespace maya::dsl;

struct Scroll {
    static constexpr int kVisible = 8;
    struct Model {
        std::vector<std::string> lines;
        int offset = 0;
    };
    struct Up {};
    struct Down {};
    struct PageUp {};
    struct PageDown {};
    struct Quit {};
    using Msg = std::variant<Up, Down, PageUp, PageDown, Quit>;

    static Model init() {
        Model m;
        for (int i = 0; i < 60; ++i) {
            m.lines.push_back(std::format("line {:03d}", i));
        }
        return m;
    }

    static auto update(Model m, Msg msg)
        -> std::pair<Model, Cmd<Msg>>
    {
        int max_off = std::max(0,
            static_cast<int>(m.lines.size()) - kVisible);
        return std::visit(overload{
            [&](Up) {
                m.offset = std::max(0, m.offset - 1);
                return std::pair{m, Cmd<Msg>{}};
            },
            [&](Down) {
                m.offset = std::min(max_off, m.offset + 1);
                return std::pair{m, Cmd<Msg>{}};
            },
            [&](PageUp) {
                m.offset = std::max(0, m.offset - kVisible);
                return std::pair{m, Cmd<Msg>{}};
            },
            [&](PageDown) {
                m.offset = std::min(max_off, m.offset + kVisible);
                return std::pair{m, Cmd<Msg>{}};
            },
            [](Quit) {
                return std::pair{Model{}, Cmd<Msg>::quit()};
            },
        }, msg);
    }

    static Element view(const Model& m) {
        std::vector<Element> visible;
        int start = m.offset;
        int end   = std::min(start + kVisible,
                             static_cast<int>(m.lines.size()));
        for (int i = start; i < end; ++i) {
            visible.push_back(text(m.lines[i]));
        }
        return v(
            v(visible),
            blank_,
            text(std::format("{} / {}",
                m.offset + 1, m.lines.size())) | Dim
        ) | pad<1> | border_<Round>;
    }

    static auto subscribe(const Model&) -> Sub<Msg> {
        return key_map<Msg>({
            {'q', Quit{}},
            {SpecialKey::Up,       Up{}},
            {SpecialKey::Down,     Down{}},
            {SpecialKey::PageUp,   PageUp{}},
            {SpecialKey::PageDown, PageDown{}},
        });
    }
};

int main() { run<Scroll>({.title = "scroll"}); }
\end{lstlisting}

\section{Notes}

\begin{itemize}[leftmargin=*]
  \item Slice in render, not in update; the Model stores
    only the offset.
  \item Clamp the offset both directions; users will mash
    arrow keys past the edges.
  \item For very long content, render a scrollbar from the
    offset and total ratio.
\end{itemize}

\section{Recap}

\begin{recap}
\begin{itemize}[leftmargin=*]
  \item Offset in Model; slice in view.
  \item Clamp at both edges.
  \item Show position as "n / total".
\end{itemize}
\end{recap}

\section*{Workshop}

\begin{enumerate}[leftmargin=*]
  \item Add a scrollbar in the right margin.
  \item Auto-scroll to the bottom when new lines arrive.
  \item Implement \code{g} to jump home, \code{G} to end.
\end{enumerate}

\chapter{Reference: \code{print} (one-shot)}
\label{ch:refman-print}

\begin{aside}
\code{maya::print} renders a single element to stdout
inline. No alt-screen, no raw mode, no event loop. The
output joins the terminal's normal scrollback and the
function returns.
\end{aside}

\section{Signature}

\begin{lstlisting}[language=C++]
namespace maya {
    void print(Element e);
}
\end{lstlisting}

\section{Working example}

\begin{lstlisting}[language=C++]
#include <maya/maya.hpp>

using namespace maya;
using namespace maya::dsl;

int main() {
    print(
        v(
            t<"build complete"> | Bold | Fg<100, 220, 120>,
            blank_,
            text("warnings: 0") | Dim,
            text("errors:   0") | Dim
        ) | pad<1> | border_<Round>
    );

    print(t<"goodbye"> | Dim);
    return 0;
}
\end{lstlisting}

\section{When to use it}

\begin{itemize}[leftmargin=*]
  \item Build tool summaries.
  \item CLI status reports.
  \item One-off cards in scripts.
  \item Anywhere you'd otherwise write a sequence of
    \code{std::cout} lines.
\end{itemize}

\section{Notes}

\begin{itemize}[leftmargin=*]
  \item Each call writes a complete frame and a newline; the
    next call starts fresh below.
  \item Style escapes are emitted only if the terminal
    supports them; piping to a file strips them.
  \item No state survives between calls.
\end{itemize}

\section{Recap}

\begin{recap}
\begin{itemize}[leftmargin=*]
  \item One element in, one frame out.
  \item Inline; no alt-screen.
  \item Perfect for CLI output.
\end{itemize}
\end{recap}

\section*{Workshop}

\begin{enumerate}[leftmargin=*]
  \item Replace a multi-line \code{printf} with \code{print}
    in an existing tool.
  \item Build a "summary card" helper that takes a title and
    a list of pairs.
  \item Pipe the output to \code{cat} and confirm the
    escapes are gone.
\end{enumerate}

\chapter{Reference: \code{live} (inline loop)}
\label{ch:refman-live}

\begin{aside}
\code{maya::live} is \code{run} without the alt-screen.
You pass a lambda that returns an Element; the runtime calls
it at the target frame rate; the rendered frame replaces
the previous one in place, in scrollback.
\end{aside}

\section{Signature}

\begin{lstlisting}[language=C++]
namespace maya {
    struct LiveConfig { int fps = 30; };
    void live(LiveConfig cfg, std::function<Element()> view);
    void quit();
}
\end{lstlisting}

\section{Working example}

\begin{lstlisting}[language=C++]
#include <maya/maya.hpp>
#include <chrono>

using namespace maya;
using namespace maya::dsl;

int main() {
    auto start = std::chrono::steady_clock::now();

    live({.fps = 30}, [&]() -> Element {
        auto elapsed =
            std::chrono::duration_cast<std::chrono::milliseconds>(
                std::chrono::steady_clock::now() - start);

        if (elapsed >= std::chrono::seconds(5)) {
            quit();
        }

        float frac = std::min(1.0,
            elapsed.count() / 5000.0);
        int filled = static_cast<int>(frac * 30);
        std::string bar;
        for (int i = 0; i < 30; ++i) {
            bar += (i < filled) ? "#" : ".";
        }
        return h(
            text(bar) | Fg<100, 220, 160>,
            text(" "),
            text(std::format("{:>3}\%",
                static_cast<int>(frac * 100))) | Bold
        );
    });

    print(t<"done"> | Bold | Fg<100, 255, 140>);
    return 0;
}
\end{lstlisting}

\section{Notes}

\begin{itemize}[leftmargin=*]
  \item Call \code{quit()} from inside the lambda to break
    the loop after this frame.
  \item Output stays in scrollback; redirecting to a file
    captures the last frame.
  \item The lambda runs on the main thread; do not block
    inside it.
\end{itemize}

\section{Recap}

\begin{recap}
\begin{itemize}[leftmargin=*]
  \item Lambda returns an Element each frame.
  \item Inline; no alt-screen.
  \item \code{quit()} ends the loop.
\end{itemize}
\end{recap}

\section*{Workshop}

\begin{enumerate}[leftmargin=*]
  \item Replace a \code{while(true) sleep} progress loop
    with \code{live}.
  \item Show a spinner during a long subprocess call.
  \item Render a small dashboard for the duration of a
    deploy.
\end{enumerate}

\chapter{Reference: focus and tab order}
\label{ch:refman-focus}

\begin{aside}
When an app has multiple focusable regions (form fields,
panels, a sidebar), one of them is "focused" at any moment.
Track that in the Model and route keys accordingly.
\end{aside}

\section{Working example}

\begin{lstlisting}[language=C++]
#include <maya/maya.hpp>

using namespace maya;
using namespace maya::dsl;

struct Focus {
    enum class Panel { Left, Right };
    struct Model {
        Panel which = Panel::Left;
        int left = 0;
        int right = 0;
    };
    struct Switch {};
    struct Bump {};
    struct Quit {};
    using Msg = std::variant<Switch, Bump, Quit>;

    static Model init() { return {}; }

    static auto update(Model m, Msg msg)
        -> std::pair<Model, Cmd<Msg>>
    {
        return std::visit(overload{
            [&](Switch) {
                m.which = (m.which == Panel::Left)
                    ? Panel::Right : Panel::Left;
                return std::pair{m, Cmd<Msg>{}};
            },
            [&](Bump) {
                if (m.which == Panel::Left) m.left++;
                else                        m.right++;
                return std::pair{m, Cmd<Msg>{}};
            },
            [](Quit) {
                return std::pair{Model{}, Cmd<Msg>::quit()};
            },
        }, msg);
    }

    static Element view(const Model& m) {
        auto panel = [&](std::string name, int v, bool active) {
            auto body = v(
                text(name) | Bold,
                blank_,
                text(std::to_string(v))
            ) | pad<1>;
            return active
                ? body | border_<Heavy>
                : body | border_<Round> | Dim;
        };

        return v(
            h(
                panel("left",  m.left,  m.which == Panel::Left)  | flex<1>,
                panel("right", m.right, m.which == Panel::Right) | flex<1>
            ),
            blank_,
            t<"tab switch, space bump, q quit"> | Dim
        );
    }

    static auto subscribe(const Model&) -> Sub<Msg> {
        return key_map<Msg>({
            {'q',             Quit{}},
            {SpecialKey::Tab, Switch{}},
            {' ',             Bump{}},
        });
    }
};

int main() { run<Focus>({.title = "focus"}); }
\end{lstlisting}

\section{Notes}

\begin{itemize}[leftmargin=*]
  \item Make the active panel visually obvious; heavy border
    plus undimmed text beats subtle marks.
  \item For many panels, store a generic index and a fixed
    array of panel handlers.
  \item Always include a visible hint for the focus key.
\end{itemize}

\section{Recap}

\begin{recap}
\begin{itemize}[leftmargin=*]
  \item One enum (or index) in the Model.
  \item Tab cycles; render distinguishes.
  \item Active panel handles the contextual keys.
\end{itemize}
\end{recap}

\section*{Workshop}

\begin{enumerate}[leftmargin=*]
  \item Add a third panel and adjust the cycle.
  \item Use \code{shift-tab} for backwards cycle.
  \item Make the panels show different content per state.
\end{enumerate}

\chapter{Reference: signals}
\label{ch:refman-signals}

\begin{aside}
OS signals (\code{SIGINT}, \code{SIGTERM}, \code{SIGWINCH},
\code{SIGUSR1}) flow into \maya{} as subscriptions. Map them
to messages; let \code{update} decide what to do.
\end{aside}

\section{Forms}

\begin{itemize}[leftmargin=*]
  \item \code{Sub<Msg>::on\_signal(SIGUSR1, msg)} — fire
    \code{msg} when \code{SIGUSR1} arrives.
  \item Common signals: \code{SIGINT} (Ctrl-C),
    \code{SIGTERM} (kill), \code{SIGUSR1},
    \code{SIGUSR2}, \code{SIGHUP}.
  \item \code{SIGWINCH} (resize) is handled separately via
    \code{on\_resize}.
\end{itemize}

\section{Working example}

\begin{lstlisting}[language=C++]
#include <maya/maya.hpp>
#include <csignal>

using namespace maya;
using namespace maya::dsl;

struct SigDemo {
    struct Model {
        int interrupts = 0;
        int reloads    = 0;
    };
    struct GotInt {};
    struct GotReload {};
    struct Quit {};
    using Msg = std::variant<GotInt, GotReload, Quit>;

    static Model init() { return {}; }

    static auto update(Model m, Msg msg)
        -> std::pair<Model, Cmd<Msg>>
    {
        return std::visit(overload{
            [&](GotInt) {
                m.interrupts++;
                return std::pair{m, Cmd<Msg>{}};
            },
            [&](GotReload) {
                m.reloads++;
                return std::pair{m, Cmd<Msg>{}};
            },
            [](Quit) {
                return std::pair{Model{}, Cmd<Msg>::quit()};
            },
        }, msg);
    }

    static Element view(const Model& m) {
        return v(
            text(std::format("interrupts: {}", m.interrupts)),
            text(std::format("reloads:    {}", m.reloads)),
            blank_,
            t<"send SIGUSR1 to reload"> | Dim,
            t<"q to quit"> | Dim
        ) | pad<1> | border_<Round>;
    }

    static auto subscribe(const Model&) -> Sub<Msg> {
        return Sub<Msg>::batch(
            key_map<Msg>({ {'q', Quit{}} }),
            Sub<Msg>::on_signal(SIGINT,  GotInt{}),
            Sub<Msg>::on_signal(SIGUSR1, GotReload{})
        );
    }
};

int main() { run<SigDemo>({.title = "signals"}); }
\end{lstlisting}

\section{Notes}

\begin{itemize}[leftmargin=*]
  \item Signals arrive between frames; ordering with key
    presses is "soon", not "instant".
  \item Catching \code{SIGINT} disables the default
    Ctrl-C-to-quit behaviour; provide your own quit path.
  \item \code{SIGTERM} should generally trigger a graceful
    shutdown command.
\end{itemize}

\section{Recap}

\begin{recap}
\begin{itemize}[leftmargin=*]
  \item Signals are subscriptions.
  \item Map signal numbers to messages.
  \item Resize has its own helper.
\end{itemize}
\end{recap}

\section*{Workshop}

\begin{enumerate}[leftmargin=*]
  \item Implement \code{SIGUSR1} → reload config.
  \item Map \code{SIGTERM} to a "shutting down" overlay
    plus a graceful quit.
  \item Send \code{SIGUSR1} from another terminal:
    \code{kill -USR1 \$(pidof yourapp)}.
\end{enumerate}

\chapter{Reference: resize handling}
\label{ch:refman-resize}

\begin{aside}
Terminals are rubber. The window resizes; \maya{} re-runs
layout; your job is to keep the Model in shape and let
\code{view} adapt.
\end{aside}

\section{Working example}

\begin{lstlisting}[language=C++]
#include <maya/maya.hpp>

using namespace maya;
using namespace maya::dsl;

struct ResizeDemo {
    struct Model { int w = 80, h = 24; };
    struct OnResize { int w; int h; };
    struct Quit {};
    using Msg = std::variant<OnResize, Quit>;

    static Model init() { return {}; }

    static auto update(Model m, Msg msg)
        -> std::pair<Model, Cmd<Msg>>
    {
        return std::visit(overload{
            [&](OnResize r) {
                m.w = r.w; m.h = r.h;
                return std::pair{m, Cmd<Msg>{}};
            },
            [](Quit) {
                return std::pair{Model{}, Cmd<Msg>::quit()};
            },
        }, msg);
    }

    static Element view(const Model& m) {
        // Adapt layout based on width.
        if (m.w < 50) {
            return v(
                text(std::format("{}x{}", m.w, m.h)) | Bold,
                t<"narrow mode"> | Dim
            ) | pad<1> | border_<Round>;
        }
        return v(
            h(
                t<"left">  | Bold,
                spacer_,
                text(std::format("{}x{}", m.w, m.h)),
                spacer_,
                t<"right"> | Bold
            ),
            blank_,
            t<"wide mode"> | Dim
        ) | pad<1> | border_<Round>;
    }

    static auto subscribe(const Model&) -> Sub<Msg> {
        return Sub<Msg>::batch(
            key_map<Msg>({ {'q', Quit{}} }),
            Sub<Msg>::on_resize([](int w, int h) {
                return OnResize{w, h};
            })
        );
    }
};

int main() { run<ResizeDemo>({.title = "resize"}); }
\end{lstlisting}

\section{Notes}

\begin{itemize}[leftmargin=*]
  \item Resize messages arrive between frames. Layout is
    always against current dimensions; you don't have to
    re-trigger anything.
  \item Use width thresholds to switch between layouts
    (narrow / wide).
  \item Test below 40 columns; expect users on phone
    terminals and tmux splits.
\end{itemize}

\section{Recap}

\begin{recap}
\begin{itemize}[leftmargin=*]
  \item Subscribe with \code{on\_resize}.
  \item Store width/height in the Model.
  \item Branch \code{view} on size for responsive layouts.
\end{itemize}
\end{recap}

\section*{Workshop}

\begin{enumerate}[leftmargin=*]
  \item Show a different layout below 60 columns.
  \item Hide secondary panels when height is small.
  \item Print a friendly "too small" placeholder under
    20x5.
\end{enumerate}

\chapter{Reference: errors and recovery}
\label{ch:refman-errors}

\begin{aside}
Things will go wrong. Files won't open. Network calls will
fail. Subprocesses will exit nonzero. Handle every failure
in \code{update}; show the user something humane.
\end{aside}

\section{Working example}

\begin{lstlisting}[language=C++]
#include <maya/maya.hpp>
#include <fstream>
#include <sstream>

using namespace maya;
using namespace maya::dsl;

struct Errors {
    struct Model {
        std::string content;
        std::string error;
    };
    struct LoadOk { std::string body; };
    struct LoadErr { std::string msg; };
    struct Load {};
    struct Quit {};
    using Msg = std::variant<Load, LoadOk, LoadErr, Quit>;

    static Model init() { return {}; }

    static auto update(Model m, Msg msg)
        -> std::pair<Model, Cmd<Msg>>
    {
        return std::visit(overload{
            [&](Load) {
                auto task = []() -> std::variant<std::string, std::string> {
                    std::ifstream f("/etc/hostname");
                    if (!f) return std::variant<std::string, std::string>{
                        std::in_place_index<1>,
                        "could not open /etc/hostname"
                    };
                    std::stringstream ss;
                    ss << f.rdbuf();
                    return std::variant<std::string, std::string>{
                        std::in_place_index<0>, ss.str()
                    };
                };
                return std::pair{m, Cmd<Msg>::task(task,
                    [](auto v) -> Msg {
                        if (v.index() == 0) return LoadOk{std::get<0>(v)};
                        return LoadErr{std::get<1>(v)};
                    })};
            },
            [&](LoadOk ok) {
                m.content = ok.body;
                m.error.clear();
                return std::pair{m, Cmd<Msg>{}};
            },
            [&](LoadErr e) {
                m.error = e.msg;
                return std::pair{m, Cmd<Msg>{}};
            },
            [](Quit) {
                return std::pair{Model{}, Cmd<Msg>::quit()};
            },
        }, msg);
    }

    static Element view(const Model& m) {
        return v(
            t<"file viewer"> | Bold,
            blank_,
            when(!m.error.empty(),
                v(
                    t<"error"> | Bold | Fg<255, 80, 80>,
                    text(m.error) | Fg<255, 180, 180>
                ),
                text(m.content)),
            blank_,
            t<"l load, q quit"> | Dim
        ) | pad<1> | border_<Round>;
    }

    static auto subscribe(const Model&) -> Sub<Msg> {
        return key_map<Msg>({
            {'q', Quit{}},
            {'l', Load{}},
        });
    }
};

int main() { run<Errors>({.title = "errors"}); }
\end{lstlisting}

\section{Notes}

\begin{itemize}[leftmargin=*]
  \item Errors are data; carry them in the Model just like
    any other field.
  \item Don't throw across the task boundary; encode failures
    as variants and dispatch.
  \item Provide a retry path; a stuck error message is
    almost as bad as a crash.
\end{itemize}

\section{Recap}

\begin{recap}
\begin{itemize}[leftmargin=*]
  \item Errors are messages.
  \item Variant of OK / err results from tasks.
  \item Render the error in the same view; offer retry.
\end{itemize}
\end{recap}

\section*{Workshop}

\begin{enumerate}[leftmargin=*]
  \item Add a retry key when an error is showing.
  \item Auto-clear errors after five seconds.
  \item Log errors to a file alongside the UI.
\end{enumerate}

\chapter{Reference: exit codes and cleanup}
\label{ch:refman-exit}

\begin{aside}
\code{run} returns an exit code. Use it. Distinguish
"user quit", "error", "interrupted", "success".
\end{aside}

\section{Working example}

\begin{lstlisting}[language=C++]
#include <maya/maya.hpp>

using namespace maya;
using namespace maya::dsl;

struct Exit {
    struct Model { int code = 0; };
    struct OkQuit {};
    struct ErrQuit {};
    using Msg = std::variant<OkQuit, ErrQuit>;

    static Model init() { return {}; }

    static auto update(Model m, Msg msg)
        -> std::pair<Model, Cmd<Msg>>
    {
        return std::visit(overload{
            [&](OkQuit) {
                m.code = 0;
                return std::pair{m, Cmd<Msg>::quit()};
            },
            [&](ErrQuit) {
                m.code = 2;
                return std::pair{m, Cmd<Msg>::quit()};
            },
        }, msg);
    }

    static Element view(const Model&) {
        return v(
            t<"choose an exit code"> | Bold,
            blank_,
            t<"q  exit 0 (ok)"> | Dim,
            t<"e  exit 2 (error)"> | Dim
        ) | pad<1> | border_<Round>;
    }

    static auto subscribe(const Model&) -> Sub<Msg> {
        return key_map<Msg>({
            {'q', OkQuit{}},
            {'e', ErrQuit{}},
        });
    }
};

int main() {
    Exit::Model final;
    // The runtime returns the exit code via run<>().
    // To propagate the Model-stored code, you can capture
    // it via a global, or return it from run() if your
    // runtime version supports model-on-quit.
    int rc = run<Exit>({.title = "exit"});
    return rc;
}
\end{lstlisting}

\section{Conventional codes}

\begin{itemize}[leftmargin=*]
  \item \code{0} — success.
  \item \code{1} — generic failure.
  \item \code{2} — bad usage.
  \item \code{130} — interrupted by Ctrl-C (\code{SIGINT}).
  \item \code{143} — killed by \code{SIGTERM}.
\end{itemize}

\section{Cleanup}

\begin{itemize}[leftmargin=*]
  \item RAII handles most of it; destructors run as
    \code{main} returns.
  \item For shared state, run cleanup inside a \code{Quit}
    handler before returning the \code{quit()} command.
  \item Don't restore the terminal yourself; \maya{} does it.
\end{itemize}

\section{Recap}

\begin{recap}
\begin{itemize}[leftmargin=*]
  \item \code{run} returns the exit code.
  \item Use conventional codes; scripts depend on them.
  \item RAII handles cleanup; trust it.
\end{itemize}
\end{recap}

\section*{Workshop}

\begin{enumerate}[leftmargin=*]
  \item Add three exit paths to your app: ok, cancel,
    error.
  \item Write a shell script that branches on the exit
    code.
  \item Verify cleanup runs after a forced \code{SIGTERM}.
\end{enumerate}

\chapter{Reference: lifecycle hooks}
\label{ch:refman-lifecycle}

\begin{aside}
Three moments in a program's life are special: start, every
update, and exit. \maya{} exposes them as commands and as
the natural shape of \code{init} and \code{update}.
\end{aside}

\section{Where each hook lives}

\begin{itemize}[leftmargin=*]
  \item \textbf{Start} — \code{init()} returns the initial
    Model. Pair it with a startup command by returning from
    \code{init} into the loop via a follow-up message.
  \item \textbf{Every update} — \code{update()} runs on
    every message; that's your "tick" if you want one.
  \item \textbf{Exit} — return \code{Cmd<Msg>::quit()}; the
    runtime calls destructors as \code{main} unwinds.
\end{itemize}

\section{Startup command pattern}

If you need to fire a command at startup, define a
\code{Boot} variant and dispatch it as the first message:

\begin{lstlisting}[language=C++]
#include <maya/maya.hpp>
#include <chrono>

using namespace maya;
using namespace maya::dsl;

struct Boot {
    struct Model { std::string status = "booting..."; };
    struct Start {};
    struct Ready {};
    struct Quit {};
    using Msg = std::variant<Start, Ready, Quit>;

    static Model init() { return {}; }

    static auto update(Model m, Msg msg)
        -> std::pair<Model, Cmd<Msg>>
    {
        return std::visit(overload{
            [&](Start) {
                return std::pair{m,
                    Cmd<Msg>::after(std::chrono::seconds(1),
                                    Ready{})};
            },
            [&](Ready) {
                m.status = "ready";
                return std::pair{m, Cmd<Msg>{}};
            },
            [](Quit) {
                return std::pair{Model{}, Cmd<Msg>::quit()};
            },
        }, msg);
    }

    static Element view(const Model& m) {
        return v(
            text(m.status) | Bold,
            blank_,
            t<"q to quit"> | Dim
        ) | pad<1> | border_<Round>;
    }

    static auto subscribe(const Model&) -> Sub<Msg> {
        return Sub<Msg>::batch(
            key_map<Msg>({ {'q', Quit{}} }),
            // Send Start once on first frame.
            Sub<Msg>::on_startup(Start{})
        );
    }
};

int main() { run<Boot>({.title = "boot"}); }
\end{lstlisting}

\section{Notes}

\begin{itemize}[leftmargin=*]
  \item Use \code{on\_startup} for one-shot init tasks;
    after the first frame, it's a no-op.
  \item Cleanup belongs in destructors. Don't try to do
    "exit work" inside the \code{Quit} handler beyond
    setting flags.
\end{itemize}

\section{Recap}

\begin{recap}
\begin{itemize}[leftmargin=*]
  \item \code{init} for startup state.
  \item \code{on\_startup} for one-shot commands.
  \item RAII for exit work.
\end{itemize}
\end{recap}

\section*{Workshop}

\begin{enumerate}[leftmargin=*]
  \item Load a config file on startup, show errors inline.
  \item Print a goodbye card after \code{run} returns.
  \item Persist Model state to disk on quit.
\end{enumerate}

\chapter{Reference: composition}
\label{ch:refman-compose}

\begin{aside}
Big apps are made of small views. The trick is to write each
view as a function and call it from the parent. State stays
in one Model; rendering decomposes.
\end{aside}

\section{Working example}

\begin{lstlisting}[language=C++]
#include <maya/maya.hpp>

using namespace maya;
using namespace maya::dsl;

struct Compose {
    struct Model {
        std::string title = "dashboard";
        int cpu = 42;
        int mem = 67;
    };
    struct Quit {};
    using Msg = std::variant<Quit>;

    static Model init() { return {}; }
    static auto update(Model m, Msg)
        -> std::pair<Model, Cmd<Msg>> {
        return {m, Cmd<Msg>::quit()};
    }

    // -- Small view helpers --
    static Element header(const Model& m) {
        return h(
            text(m.title) | Bold,
            spacer_,
            t<"q to quit"> | Dim
        );
    }

    static Element gauge(std::string label, int pct) {
        std::string bar;
        for (int i = 0; i < 20; ++i) {
            bar += (i < pct / 5) ? "#" : ".";
        }
        return h(
            text(label) | Bold,
            text(": "),
            text(bar) | Fg<100, 220, 160>,
            text(std::format(" {}\%", pct))
        );
    }

    static Element body(const Model& m) {
        return v(
            gauge("cpu", m.cpu),
            gauge("mem", m.mem)
        );
    }

    // -- Top-level view glues them together --
    static Element view(const Model& m) {
        return v(
            header(m),
            blank_,
            body(m)
        ) | pad<1> | border_<Round>;
    }

    static auto subscribe(const Model&) -> Sub<Msg> {
        return key_map<Msg>({ {'q', Quit{}} });
    }
};

int main() { run<Compose>({.title = "compose"}); }
\end{lstlisting}

\section{Notes}

\begin{itemize}[leftmargin=*]
  \item Pass the Model by const reference into helpers;
    don't reach for globals.
  \item Helpers that take just the data they need are
    easier to test and refactor.
  \item For very large apps, split helpers across files; the
    \code{Element} return type travels fine across TUs.
\end{itemize}

\section{Recap}

\begin{recap}
\begin{itemize}[leftmargin=*]
  \item Helpers return \code{Element}.
  \item Pass only what they need.
  \item Top-level \code{view} glues them.
\end{itemize}
\end{recap}

\section*{Workshop}

\begin{enumerate}[leftmargin=*]
  \item Extract three helpers from your current view.
  \item Move them to a separate \code{views.cpp}.
  \item Pass a parameter to make one helper reusable.
\end{enumerate}

\chapter{Reference: testing programs}
\label{ch:refman-test}

\begin{aside}
The whole point of the Model–Update–View shape is that
testing is trivial. \code{update} is a pure function from
(Model, Msg) to (Model, Cmd). You don't need a terminal,
a fixture, or a mock; you need a unit test.
\end{aside}

\section{Working example}

\begin{lstlisting}[language=C++]
// counter.hpp
#pragma once
#include <maya/maya.hpp>

struct Counter {
    struct Model { int count = 0; };
    struct Inc {};
    struct Dec {};
    struct Quit {};
    using Msg = std::variant<Inc, Dec, Quit>;

    static Model init() { return {}; }

    static auto update(Model m, Msg msg)
        -> std::pair<Model, maya::Cmd<Msg>>
    {
        return std::visit(maya::overload{
            [&](Inc) { m.count++; return std::pair{m, maya::Cmd<Msg>{}}; },
            [&](Dec) { m.count--; return std::pair{m, maya::Cmd<Msg>{}}; },
            [](Quit) {
                return std::pair{Model{}, maya::Cmd<Msg>::quit()};
            },
        }, msg);
    }

    static maya::Element view(const Model& m) {
        return maya::dsl::text(m.count);
    }

    static auto subscribe(const Model&) -> maya::Sub<Msg> {
        return maya::dsl::key_map<Msg>({});
    }
};
\end{lstlisting}

\begin{lstlisting}[language=C++]
// counter_test.cpp
#include "counter.hpp"
#include <cassert>

int main() {
    auto m0 = Counter::init();
    assert(m0.count == 0);

    auto [m1, _1] = Counter::update(m0, Counter::Inc{});
    assert(m1.count == 1);

    auto [m2, _2] = Counter::update(m1, Counter::Inc{});
    assert(m2.count == 2);

    auto [m3, _3] = Counter::update(m2, Counter::Dec{});
    assert(m3.count == 1);

    return 0;
}
\end{lstlisting}

\section{What to test}

\begin{itemize}[leftmargin=*]
  \item \code{update} for every \code{Msg} variant.
  \item Invariants in the Model (counter never negative,
    selection always in range).
  \item Returned commands when relevant — assert
    \code{cmd.is\_quit()} on the Quit branch, for example.
\end{itemize}

\section{What \emph{not} to test}

\begin{itemize}[leftmargin=*]
  \item \code{view} pixel-by-pixel. Snapshot tests grow
    brittle and tell you nothing useful.
  \item The runtime itself. \maya{} has its own tests; trust
    them.
\end{itemize}

\section{Recap}

\begin{recap}
\begin{itemize}[leftmargin=*]
  \item Pure functions invite unit tests.
  \item Test \code{update}; assert on Model after each Msg.
  \item Skip the view.
\end{itemize}
\end{recap}

\section*{Workshop}

\begin{enumerate}[leftmargin=*]
  \item Write a property test: \code{Inc} then \code{Dec}
    is a no-op.
  \item Test that \code{Quit} always returns the quit
    command.
  \item Add a test for an edge case in your real app.
\end{enumerate}

\chapter{Reference: CMake integration}
\label{ch:refman-cmake}

\begin{aside}
\maya{} ships a CMake package config. Find it; link to
\code{maya::maya}; you're done.
\end{aside}

\section{Minimum project}

A working project has three files: \code{CMakeLists.txt},
\code{main.cpp}, and a build directory.

\begin{lstlisting}[language=cmake]
# CMakeLists.txt
cmake_minimum_required(VERSION 3.20)
project(hello CXX)

set(CMAKE_CXX_STANDARD 23)
set(CMAKE_CXX_STANDARD_REQUIRED ON)

find_package(maya REQUIRED)

add_executable(hello main.cpp)
target_link_libraries(hello PRIVATE maya::maya)
\end{lstlisting}

\begin{lstlisting}[language=C++]
// main.cpp
#include <maya/maya.hpp>

using namespace maya;
using namespace maya::dsl;

int main() {
    print(t<"hello"> | Bold | Fg<100, 220, 160>);
}
\end{lstlisting}

\section{Build it}

\begin{lstlisting}[language=bash]
cmake -S . -B build
cmake --build build -j
./build/hello
\end{lstlisting}

\section{Linking against widgets}

Widgets live in headers; they pull no extra targets. Just
include them:

\begin{lstlisting}[language=C++]
#include <maya/widget/inline_diff.hpp>
\end{lstlisting}

\section{Vendoring instead of installing}

If you want \maya{} as a submodule:

\begin{lstlisting}[language=cmake]
add_subdirectory(third_party/maya)
target_link_libraries(hello PRIVATE maya::maya)
\end{lstlisting}

\section{Notes}

\begin{itemize}[leftmargin=*]
  \item \code{find\_package(maya REQUIRED)} fails loud if
    the package isn't installed; that's good.
  \item Use \code{target\_compile\_features(... cxx\_std\_23)}
    for older CMake versions that don't honour the global
    \code{CMAKE\_CXX\_STANDARD}.
  \item For sanitisers, add \code{-fsanitize=address,undefined}
    to \code{target\_compile\_options} in debug builds.
\end{itemize}

\section{Recap}

\begin{recap}
\begin{itemize}[leftmargin=*]
  \item \code{find\_package(maya REQUIRED)}.
  \item Link \code{maya::maya}.
  \item Widgets are header-only.
\end{itemize}
\end{recap}

\section*{Workshop}

\begin{enumerate}[leftmargin=*]
  \item Create a fresh \maya{} project from scratch in five
    minutes.
  \item Add a unit test target via \code{add\_test}.
  \item Wire up a sanitiser-only build configuration.
\end{enumerate}

\chapter{Reference: complete example apps}
\label{ch:refman-apps}

\begin{aside}
The previous chapters showed one feature at a time. This
one shows three short but complete programs. Each is small
enough to read in one sitting and exercises several
features at once.
\end{aside}

\section{App 1: ToDo list}

\begin{lstlisting}[language=C++]
#include <maya/maya.hpp>
#include <vector>
#include <string>

using namespace maya;
using namespace maya::dsl;

struct Todo {
    struct Item { std::string text; bool done = false; };
    struct Model {
        std::vector<Item> items;
        std::string buffer;
        int selected = 0;
    };
    struct Type { char c; };
    struct Back {};
    struct Add {};
    struct Toggle {};
    struct Up {};
    struct Down {};
    struct Quit {};
    using Msg = std::variant<Type, Back, Add, Toggle, Up, Down, Quit>;

    static Model init() { return {}; }

    static auto update(Model m, Msg msg)
        -> std::pair<Model, Cmd<Msg>>
    {
        return std::visit(overload{
            [&](Type t) { m.buffer.push_back(t.c); return std::pair{m, Cmd<Msg>{}}; },
            [&](Back) {
                if (!m.buffer.empty()) m.buffer.pop_back();
                return std::pair{m, Cmd<Msg>{}};
            },
            [&](Add) {
                if (!m.buffer.empty()) {
                    m.items.push_back({m.buffer, false});
                    m.buffer.clear();
                }
                return std::pair{m, Cmd<Msg>{}};
            },
            [&](Toggle) {
                if (!m.items.empty() && m.selected < static_cast<int>(m.items.size())) {
                    m.items[m.selected].done = !m.items[m.selected].done;
                }
                return std::pair{m, Cmd<Msg>{}};
            },
            [&](Up) {
                m.selected = std::max(0, m.selected - 1);
                return std::pair{m, Cmd<Msg>{}};
            },
            [&](Down) {
                int n = static_cast<int>(m.items.size());
                m.selected = std::min(std::max(0, n - 1), m.selected + 1);
                return std::pair{m, Cmd<Msg>{}};
            },
            [](Quit) { return std::pair{Model{}, Cmd<Msg>::quit()}; },
        }, msg);
    }

    static Element view(const Model& m) {
        std::vector<Element> rows;
        for (size_t i = 0; i < m.items.size(); ++i) {
            auto& it = m.items[i];
            auto mark = it.done ? "[x] " : "[ ] ";
            auto line = h(text(mark), text(it.text));
            if (static_cast<int>(i) == m.selected) line = line | Bold | Bg<40, 80, 140>;
            if (it.done) line = line | Dim;
            rows.push_back(line);
        }
        return v(
            t<"todo"> | Bold,
            blank_,
            v(rows),
            blank_,
            h(t<"add: ">, text(m.buffer), text("_") | Blink),
            blank_,
            t<"enter add, space toggle, up/down move, q quit"> | Dim
        ) | pad<1> | border_<Round>;
    }

    static auto subscribe(const Model&) -> Sub<Msg> {
        std::vector<std::pair<KeyEvent, Msg>> binds;
        for (char c = ' '; c < 127; ++c) binds.push_back({c, Type{c}});
        binds.push_back({SpecialKey::Backspace, Back{}});
        binds.push_back({SpecialKey::Enter,     Add{}});
        binds.push_back({' ',                   Toggle{}});
        binds.push_back({SpecialKey::Up,        Up{}});
        binds.push_back({SpecialKey::Down,      Down{}});
        binds.push_back({ctrl('c'),             Quit{}});
        return key_map<Msg>(std::move(binds));
    }
};

int main() { run<Todo>({.title = "todo"}); }
\end{lstlisting}

\section{App 2: live system stats}

A version of \code{sysmon} small enough to fit on one page,
suitable as a starting point.

\begin{lstlisting}[language=C++]
#include <maya/maya.hpp>
#include <chrono>
#include <random>

using namespace maya;
using namespace maya::dsl;

struct Stats {
    struct Model {
        std::vector<float> cpu;
        std::vector<float> mem;
        std::mt19937 rng{42};
    };
    struct Tick {};
    struct Quit {};
    using Msg = std::variant<Tick, Quit>;

    static Model init() { return {}; }

    static auto update(Model m, Msg msg)
        -> std::pair<Model, Cmd<Msg>>
    {
        return std::visit(overload{
            [&](Tick) {
                std::uniform_real_distribution<float> d(0.f, 1.f);
                m.cpu.push_back(d(m.rng));
                m.mem.push_back(d(m.rng));
                if (m.cpu.size() > 40) m.cpu.erase(m.cpu.begin());
                if (m.mem.size() > 40) m.mem.erase(m.mem.begin());
                return std::pair{m, Cmd<Msg>{}};
            },
            [](Quit) { return std::pair{Model{}, Cmd<Msg>::quit()}; },
        }, msg);
    }

    static Element spark(const std::vector<float>& data) {
        const char* glyphs[] = {".", ":", "-", "=", "*", "#"};
        std::string s;
        for (float v : data) {
            int idx = std::min(5,
                std::max(0, static_cast<int>(v * 6)));
            s += glyphs[idx];
        }
        return text(s);
    }

    static Element view(const Model& m) {
        return v(
            t<"system">  | Bold,
            blank_,
            h(t<"cpu: ">, spark(m.cpu) | Fg<255, 200, 100>),
            h(t<"mem: ">, spark(m.mem) | Fg<100, 200, 255>),
            blank_,
            t<"q quit"> | Dim
        ) | pad<1> | border_<Round>;
    }

    static auto subscribe(const Model&) -> Sub<Msg> {
        return Sub<Msg>::batch(
            key_map<Msg>({ {'q', Quit{}} }),
            Sub<Msg>::every(std::chrono::milliseconds(200),
                            Tick{})
        );
    }
};

int main() { run<Stats>({.title = "stats"}); }
\end{lstlisting}

\section{Notes}

\begin{itemize}[leftmargin=*]
  \item Both apps fit in under 100 lines.
  \item Both are pure aside from the timer (and the
    todo's buffer-as-state).
  \item Both can be extended without restructuring.
\end{itemize}

\section{Recap}

\begin{recap}
\begin{itemize}[leftmargin=*]
  \item Small apps are a few dozen lines each.
  \item Same skeleton every time.
  \item The grammar is small; the vocabulary is wide.
\end{itemize}
\end{recap}

\section*{Workshop}

\begin{enumerate}[leftmargin=*]
  \item Add a "clear all" key to the todo.
  \item Replace the random source in stats with real \code{/proc} reads.
  \item Combine both into a tab-switched app.
\end{enumerate}

\chapter{Reference: snippets cookbook}
\label{ch:refman-cookbook}

\begin{aside}
Small, copy-pasteable patterns for the situations that come
up every day. Each is a complete fragment; treat them as
slot-ins.
\end{aside}

\section{Conditional border}

\begin{lstlisting}[language=C++]
auto card = body | pad<1>;
return active ? card | border_<Heavy> : card | border_<Round>;
\end{lstlisting}

\section{Build a row from a vector}

\begin{lstlisting}[language=C++]
std::vector<Element> children;
for (auto& item : items) children.push_back(text(item));
return h(children);
\end{lstlisting}

\section{Dim everything except the focused row}

\begin{lstlisting}[language=C++]
for (size_t i = 0; i < rows.size(); ++i) {
    if (static_cast<int>(i) != selected) rows[i] = rows[i] | Dim;
}
\end{lstlisting}

\section{Right-justify a value}

\begin{lstlisting}[language=C++]
h(text(label), spacer_, text(value) | Bold)
\end{lstlisting}

\section{Centred title above content}

\begin{lstlisting}[language=C++]
v(
    center(text(title) | Bold),
    blank_,
    content
)
\end{lstlisting}

\section{Status bar at the bottom}

\begin{lstlisting}[language=C++]
return v(
    body | flex<1>,
    h(t<"ready">, spacer_, text(time)) | Bg<30, 30, 40>
);
\end{lstlisting}

\section{Quit on Escape or Ctrl-C}

\begin{lstlisting}[language=C++]
return key_map<Msg>({
    {SpecialKey::Escape, Quit{}},
    {ctrl('c'),          Quit{}},
});
\end{lstlisting}

\section{Toggle a flag with a key}

\begin{lstlisting}[language=C++]
[&](Toggle) { m.flag = !m.flag; return std::pair{m, Cmd<Msg>{}}; }
\end{lstlisting}

\section{Auto-clear flash after a delay}

\begin{lstlisting}[language=C++]
[&](Flash) {
    m.show = true;
    return std::pair{m,
        Cmd<Msg>::after(std::chrono::milliseconds(500),
                        ClearFlash{})};
}
\end{lstlisting}

\section{Two-column responsive layout}

\begin{lstlisting}[language=C++]
if (width >= 80) {
    return h(left | flex<1>, right | flex<1>);
}
return v(left, right);
\end{lstlisting}

\section{Show a "loading" placeholder}

\begin{lstlisting}[language=C++]
when(m.loading,
    t<"loading..."> | Dim,
    text(m.data));
\end{lstlisting}

\section{Number formatting helpers}

\begin{lstlisting}[language=C++]
auto pct  = std::format("{:>3}\%", static_cast<int>(f * 100));
auto kb   = std::format("{:>6} KB", bytes / 1024);
auto time = std::format("{:02}:{:02}", mins, secs);
\end{lstlisting}

\section{A horizontal rule}

\begin{lstlisting}[language=C++]
text(std::string(40, '-')) | Dim
\end{lstlisting}

\section{Notes}

\begin{itemize}[leftmargin=*]
  \item Snippets above are fragments; drop them into a real
    \code{Program} skeleton.
  \item Adapt the literals; the shape is what matters.
\end{itemize}

\section{Recap}

\begin{recap}
\begin{itemize}[leftmargin=*]
  \item A dozen patterns; every day's UI work.
  \item Copy, paste, adapt.
\end{itemize}
\end{recap}

\section*{Workshop}

\begin{enumerate}[leftmargin=*]
  \item Add three patterns of your own from your last
    project.
  \item Wrap each as a function returning \code{Element}.
  \item Build a "starter kit" header.
\end{enumerate}

\chapter{Reference: full enum and key list}
\label{ch:refman-enums}

\begin{aside}
The complete set of enumerated constants \maya{} exposes.
Look here when the compiler asks for a name and you've
forgotten it.
\end{aside}

\section{SpecialKey}

\begin{itemize}[leftmargin=*]
  \item \code{Up}, \code{Down}, \code{Left}, \code{Right}.
  \item \code{Enter}, \code{Escape}, \code{Tab},
    \code{Backspace}, \code{Delete}.
  \item \code{Home}, \code{End}, \code{PageUp},
    \code{PageDown}, \code{Insert}.
  \item \code{F1} through \code{F12}.
\end{itemize}

\section{BorderStyle}

\begin{itemize}[leftmargin=*]
  \item \code{None}, \code{Round}, \code{Square},
    \code{Heavy}, \code{Double}, \code{Ascii}.
\end{itemize}

\section{MouseButton}

\begin{itemize}[leftmargin=*]
  \item \code{Left}, \code{Middle}, \code{Right}.
  \item \code{WheelUp}, \code{WheelDown},
    \code{WheelLeft}, \code{WheelRight}.
  \item \code{None} for move events with no button held.
\end{itemize}

\section{MouseKind}

\begin{itemize}[leftmargin=*]
  \item \code{Press}, \code{Release}, \code{Drag},
    \code{Move}.
\end{itemize}

\section{Cursor}

\begin{itemize}[leftmargin=*]
  \item \code{Hidden}, \code{Block}, \code{Underscore},
    \code{Bar}.
  \item \code{BlinkingBlock}, \code{BlinkingBar},
    \code{BlinkingUnderscore}.
\end{itemize}

\section{Working example}

\begin{lstlisting}[language=C++]
#include <maya/maya.hpp>

using namespace maya;
using namespace maya::dsl;

struct EnumDemo {
    struct Model { std::string last = "(none)"; };
    struct Got { std::string label; };
    struct Quit {};
    using Msg = std::variant<Got, Quit>;

    static Model init() { return {}; }
    static auto update(Model m, Msg msg)
        -> std::pair<Model, Cmd<Msg>> {
        return std::visit(overload{
            [&](Got g) { m.last = g.label; return std::pair{m, Cmd<Msg>{}}; },
            [](Quit) { return std::pair{Model{}, Cmd<Msg>::quit()}; },
        }, msg);
    }

    static Element view(const Model& m) {
        return v(
            text(std::format("last: {}", m.last)),
            blank_,
            t<"try F1..F4, arrows, page up/down"> | Dim,
            t<"q to quit"> | Dim
        ) | pad<1> | border_<Round>;
    }

    static auto subscribe(const Model&) -> Sub<Msg> {
        return key_map<Msg>({
            {'q', Quit{}},
            {SpecialKey::F1,       Got{"F1"}},
            {SpecialKey::F2,       Got{"F2"}},
            {SpecialKey::F3,       Got{"F3"}},
            {SpecialKey::F4,       Got{"F4"}},
            {SpecialKey::Up,       Got{"up"}},
            {SpecialKey::Down,     Got{"down"}},
            {SpecialKey::Left,     Got{"left"}},
            {SpecialKey::Right,    Got{"right"}},
            {SpecialKey::PageUp,   Got{"pgup"}},
            {SpecialKey::PageDown, Got{"pgdn"}},
            {SpecialKey::Home,     Got{"home"}},
            {SpecialKey::End,      Got{"end"}},
        });
    }
};

int main() { run<EnumDemo>({.title = "enums"}); }
\end{lstlisting}

\section{Notes}

\begin{itemize}[leftmargin=*]
  \item Not every key reaches every terminal; F-keys in
    particular are sometimes captured by tmux or by the
    terminal emulator itself.
  \item \code{Insert} is rarely useful; default to other
    bindings.
  \item \code{shift(...)}, \code{ctrl(...)}, \code{alt(...)}
    can modify any of these.
\end{itemize}

\section{Recap}

\begin{recap}
\begin{itemize}[leftmargin=*]
  \item One enum per concern.
  \item All keys reachable via \code{key\_map}.
  \item Modifier wrappers compose.
\end{itemize}
\end{recap}

\section*{Workshop}

\begin{enumerate}[leftmargin=*]
  \item Print a label for every key you press.
  \item Build a "vim-like" set of bindings (\code{hjkl},
    \code{gg}, \code{G}).
  \item Try every cursor style with \code{RunConfig::cursor}.
\end{enumerate}

\include{chapters/13-events}
\include{chapters/14-cmdsub}
\include{chapters/15-widgets}
\include{chapters/16-build}

\backmatter
\include{chapters/99-colophon}

\end{document}
